%Basic constructions like size and language :
\documentclass[12pt]{report}
\usepackage[utf8]{inputenc}
\usepackage{CJKutf8} %Chinese, support of begin cjk end cjk
\usepackage[T1]{fontenc} %French
\usepackage[letterpaper, margin=1.27cm]{geometry}
\setlength{\emergencystretch}{3em} %the very wide text block needs slack to avoid overfull lines
\usepackage{subfiles} 
\usepackage{times} %Times New Roman font

%References:
\usepackage[round]{natbib} %author-year citations, loaded before hyperref

%Make title, table of contents
\usepackage{hyperref} %necessay for maketitle to work
\hypersetup{colorlinks=true, allcolors=blue} %all links and citations in blue
\usepackage{datetime2} %for the dtmnow command

\title{Casual Document for Probabilistic Principal Component Analysis MLE Asymptotic Normality}
\author{Wen, Zehai}
\date{\DTMnow}

%Page style, header and footer:
\usepackage{fancyhdr}

\pagestyle{fancy}
\setlength{\headheight}{14.5pt}
\addtolength{\topmargin}{-2.5pt}
\fancyhf{}
\lhead{
Wen, Zehai}
\rhead{\today}

\rfoot{}
\cfoot{\thepage}

%Graphics and figures:
\usepackage{graphicx} 
\usepackage{float} %both graphicx and float packages are necessary for begin figure
\usepackage{caption} %necessary for any kind of caption
\usepackage{subcaption} %necessary for subfigure

%Tabular:
\usepackage{multirow}

%Coding:
\usepackage{listings}
\usepackage{color}

\lstset{literate=%
  {é}{{\'e}}{1}%
  {è}{{\`e}}{1}%
  {à}{{\`a}}{1}%
  {ç}{{\c{c}}}{1}%
  {œ}{{\oe}}{1}%
  {ù}{{\`u}}{1}%
  {É}{{\'E}}{1}%
  {È}{{\`E}}{1}%
  {À}{{\`A}}{1}%
  {Ç}{{\c{C}}}{1}%
  {Œ}{{\OE}}{1}%
  {Ê}{{\^E}}{1}%
  {ê}{{\^e}}{1}%
  {î}{{\^i}}{1}%
  {ô}{{\^o}}{1}%
  {û}{{\^u}}{1}%
  {ä}{{\"{a}}}1
  {ë}{{\"{e}}}1
  {ï}{{\"{i}}}1
  {ö}{{\"{o}}}1
  {ü}{{\"{u}}}1
  {û}{{\^{u}}}1
  {â}{{\^{a}}}1
  {Â}{{\^{A}}}1
  {Î}{{\^{I}}}1
} 

\definecolor{dkgreen}{rgb}{0,0.6,0}
\definecolor{gray}{rgb}{0.5,0.5,0.5}
\definecolor{mauve}{rgb}{0.58,0,0.82}

\lstset{frame=tb,
  language=Java,
  aboveskip=3mm,
  belowskip=3mm,
  showstringspaces=false,
  columns=flexible,
  basicstyle={\small\ttfamily},
  numbers=none,
  numberstyle=\tiny\color{gray},
  keywordstyle=\color{blue},
  commentstyle=\color{dkgreen},
  stringstyle=\color{mauve},
  breaklines=true,
  breakatwhitespace=true,
  tabsize=3
}

%Graph Theory:
\usepackage{tikz}
\usetikzlibrary{arrows.meta}

%Commands :
\newcommand\tab[1][1cm]{\hspace*{#1}}
\renewcommand{\bibname}{Références}

%%%%%%%%%%%%%%%%%%%%%%%%%%%%Mathematical packages%%%%%%%%%%%%%%%%%%%%%%%%%%%%%%%%%%%%%%%%%%

%Basic utilities:
\usepackage{amssymb}
\usepackage{amsmath}
\usepackage{amsfonts}
\usepackage{physics} %for vector symbol
\usepackage{bm} %for bm

%Commands :
\DeclareMathOperator*{\argmax}{arg\,max}
\DeclareMathOperator*{\argmin}{arg\,min}

\newcommand{\bracket}[1]{\left ( #1 \right )}
\newcommand{\curly}[1]{\left \{ #1 \right \} }
\newcommand{\squared}[1]{\left [ #1 \right ]}

\usepackage{amsthm}  %Theorem styles
\theoremstyle{plain}
\newtheorem{theorem}{Theorem}

\theoremstyle{definition}
\newtheorem{definition}[theorem]{Definition} 
\newtheorem{definitions}[theorem]{Definitions} 
\newtheorem{example}[theorem]{Example}
\newtheorem{examples}[theorem]{Examples}
\newtheorem{remark}[theorem]{Remark}
\newtheorem{remarks}[theorem]{Remarks}
\newtheorem{lemma}[theorem]{Lemma}
\newtheorem{corollary}[theorem]{Corollary}
\newtheorem{proposition}[theorem]{Proposition}

%%%%%%%%%%%%%%%%%%%%%%%%%%%%%%%%%%%%%%%%%%%%%%%%%%%%%%%%%%%%%%%%%%%%%%%%%%%%%%%%

\begin{document}
\begin{CJK*}{UTF8}{gbsn}
\maketitle

\setcounter{tocdepth}{0}
%\tableofcontents

\chapter{Low Dimensional Asymptotic Normality}

\section{Basic Setup}\label{sec:basic_setup}

Given a "high dimensional" data set $\{\bm{x}_i\}_{i=1}^n$
where $\bm{x}_i \in \mathbb{R}^p$ for all $i = 1, \ldots, n$,
the statistician fixes a positive integer $q < p$ 
and wishes to find the $q$ principal components.
Assuming that the data is centred (that is, with mean $\bm{0}$),
the probabilistic principal component analysis model
is a statistical model that
treats $\bm{x}_i$'s as realizations 
of independent random variables $\bm{X}_i = \bm{W}\bm{Z}_i + \bm{\epsilon_i}$,
where $\bm{W}$ is an unknown $p \times q$ loading matrix of full rank, 
$\bm{Z}_i \sim N_q\bracket{\bm{0}, \bm{I}_q}$ is the latent normal variable,
and $\bm{\epsilon_i} \sim N_p\bracket{\bm{0}, \sigma^2\bm{I}_p}$
is the Gaussian noise with unknown standard deviation $\sigma > 0$,
and $(\bm{Z}_1, \ldots, \bm{Z}_n, \bm{\epsilon}_1, \ldots, \bm{\epsilon}_n)$
are all independent (for now). 
Put $\theta = (\bm{W}_{\theta}, \sigma_{\theta}^2)$ for the unknown 
parameter to estimate
and $\Theta$ for all possible choices of $\theta$.
$\Theta$ is thus the collection of all tuples $(\bm{W}, \sigma^2)$
such that $\bm{W}$ is a real $p \times q$ matrix with $\text{rank}(\bm{W}) = q$
and $\sigma^2 \in (0, \infty)$.
Given that the collection of all real $p \times q$ matrix with rank $q$
is an open subset of the Euclidean space $\mathbb{R}^{pq}$,
we consider $\Theta$ as an open subset of the Euclidean space $\mathbb{R}^{pq+1}$
and our statistical model is 

\begin{equation}\label{eq:original_model}
  \mathcal{P} = \curly{N_p\bracket{\bm{0}, \bm{W}_{\theta}\bm{W}_{\theta}^T + \sigma_{\theta}^2\bm{I}_p}}_{\theta \in \Theta}.
\end{equation}

To resolve the non-identifiability of $\mathcal{P}$,
we introduce the quotient manifold structure.
For $(\bm{A},s) \in \Theta$ and $(\bm{B}, t) \in \Theta$,
we define $(\bm{A},s) \sim (\bm{B}, t)$
to mean that they give rise to the same covariance matrix, namely when
$\bm{A}\bm{A}^T + s\bm{I}_p = \bm{B}\bm{B}^T + t\bm{I}_p$.
The equivalence classes are denoted by $[(\bm{A},s)]$
or more briefly $[\bm{A},s]$. 
Unless otherwise mentioned, $(\bm{A},s) \in \Theta$ is the default representative 
we pick when we write an equivalence class as $[\bm{A},s]$.
Writing $O(q)$ for the collection of all $q \times q$ real orthogonal matrices,
it was shown in Datta's thesis that 
$[\bm{A},s] = \{(\bm{A}\bm{Q},s) \mid \bm{Q} \in O(q)\}$.
Therefore, $M := \Theta / \sim$ can be identified as the orbits 
of the right action of $O(q)$ on $\Theta$
and, given that this Lie group action is smooth, proper, and free,
the quotient manifold theorem (see \citealp{Lee_2013}, Theorem 21.10) shows that there is a unique smooth 
structure making $M$ a manifold of dimension 

\begin{equation}\label{eq:d_def}
  d = pq+1-\frac{q(q-1)}{2}
\end{equation} 
and 
the map $\pi: \Theta \to M$ 
defined by

\begin{equation}\label{eq:pi_def}
  \pi((\bm{A},s)) = [\bm{A},s]
\end{equation}
is a smooth submersion.
Moreover, this action is vertical, transitive on fibers 
and is isometric with respect to the Euclidean metric.
Thus, there exists a unique metric $g$ on $M$ making
$(M,g)$ a Riemannian manifold and $\pi$ a Riemannian submersion
(see \citealp{Lee_2018}, Theorem 2.28).
In more concrete terms, 
if $g^{\Theta}$ is the standard Euclidean metric given by 
\begin{equation}\label{eq:euclidean_metric}
  g_{(\bm{A},s)}^{\Theta}\bracket{(\bm{V},c), (\bm{U},b)} = \tr(\bm{V}\bm{U}^T) + cb,
\end{equation}
whenever $(\bm{A},s) \in \Theta$, $(\bm{V},c)$, 
$(\bm{U},b) \in T_{(\bm{A},s)}\Theta \cong \mathbb{R}^{p \times q} \times \mathbb{R}$,
then 
\begin{equation}\label{eq:g_def}
  g_{[\bm{A},s]}(\bm{v}, \bm{h})
  = g_{[\bm{A},s]} \bracket{d\pi_{(A,s)}((\bm{V},c)), d\pi_{(A,s)}((\bm{U},b))}
  = g_{(\bm{A},s)}^{\Theta}\bracket{(\bm{V},c), (\bm{U},b)},
\end{equation}
whenever $\bm{v}$, $\bm{h} \in T_{[\bm{A},s]}M$
lifting to $(\bm{V},c)$ and $(\bm{U},b)$ respectively.
A standard fact is that $\Theta$ is connected.
Consequently, $M$ is also connected.
For later convenience,
let us define two functions $\Phi: \Theta \to \mathbb{R}^{p \times p}$
and $\bm{\Sigma}: M \to \mathbb{R}^{p \times p}$
by setting, for every $[\bm{A},s] \in M$:

\begin{equation}\label{eq:Sigma_Phi_def}
  \bm{\Sigma}([\bm{A},s]) := \Phi((\bm{A},s)) := \bm{A}\bm{A}^T + s\bm{I}_p.
\end{equation}
In other words, $\bm{\Sigma} \circ \pi = \Phi$
gives the covariance matrix.
\eqref{eq:original_model} is then reparametrized as

\begin{equation}\label{eq:quotient_model}
  \mathcal{P} = \curly{N_p\bracket{\bm{0}, \bm{\Sigma}([\bm{A},s])}}_{[\bm{A},s]\in M}
\end{equation}
and this is the identifiable model we shall work with.

Finally, thoughout the entire document,
our default norm for a matrix $\bm{A}$
will be the spectral norm, denoted by $\|\bm{A}\|$.
We will also make use of the 
Frobenius norm $\|\bm{A}\|_F$ as well.
The eigenvalues and the singular values of a matrix 
will always be labelled in decreasing order.
In symbols, if $\bm{A}$ is a matrix of size $p \times q$,
then we denote its $q$ singular values as 

\begin{equation*}
  \|\bm{A}\| = \sigma_1(\bm{A}) \geqslant \sigma_2(\bm{A}) \geqslant \cdots \geqslant 
  \sigma_q(\bm{A}) = \sigma_{\min}(\bm{A}) \geqslant 0
\end{equation*}
and, if $\bm{S}$ is a symmetric matrix of size $p \times p$,
then we denote its $p$ eigenvalues as 

\begin{equation*}
  \lambda_1(\bm{S}) \geqslant \lambda_2(\bm{S}) 
  \geqslant \cdots \geqslant \lambda_p(\bm{S}).
\end{equation*}

\section{Geometric Facts about $(M,g)$}

\subsection{Horizontal and Verticle Tangent Spaces}
Working with Riemannian submersion, 
we shall as usual start with computations regarding the horizontal and vertical 
tangent spaces.

\begin{proposition}\label{thm:horizontal_vertical}
\,
\begin{enumerate}
  \item The map $\bm{\Sigma}$ defined in \eqref{eq:Sigma_Phi_def}
is a smooth injective immersion.
  \item At every $(\bm{A},s) \in \Theta$,
  the vertical tangent space 

  \begin{equation}\label{eq:vertical_space}
  V_{(\bm{A},s)} :=\ker(\mathrm{d}\pi_{(\bm{A},s)}) = \ker(\mathrm{d}\Phi_{(\bm{A},s)})
  \end{equation}
  equals to the collection of all tangent vectors 
of the form $(\bm{A} \bm{\Omega}, 0) \in T_{(\bm{A},s)}\Theta$,
with $\bm{\Omega}$ ranged over all $q \times q$ skew symmetric matrices.
  \item At every $(\bm{A},s) \in \Theta$,
  the horizontal tangent space $H_{(\bm{A},s)} := V_{(\bm{A},s)}^{\perp}$
equals to the collection of all tangent vectors 
of the form $(\bm{V}, c) \in T_{(\bm{A},s)}\Theta$
ranged over all $c \in \mathbb{R}$ and all 
$p \times q$ matrices $\bm{V}$
satisfying $\bm{V}^T \bm{A} = \bm{A}^T \bm{V}$.
\end{enumerate}
\end{proposition}

\begin{proof}
Fix a point $(\bm{A},s) \in \Theta$
throughout this proof.
Before we start the proof,
let us first study the set $\ker(\mathrm{d}\Phi_{(\bm{A},s)})$,
where $\Phi$ is given in \eqref{eq:Sigma_Phi_def}.
For $(\bm{V}, c) \in T_{(\bm{A},s)}\Theta$, 
define a curve $\gamma : (-\epsilon, \epsilon) \to \Theta$
by setting $\gamma(t) = (\bm{A}+t\bm{V}, s + tc)$.
Then:

\begin{equation*}
  \mathrm{d} \Phi_{(\bm{A},s)} ((\bm{V}, c)) = 
  \frac{d}{dt} \bigg|_{t = 0} \Phi(\gamma(t)) = \bm{A}\bm{V}^T + \bm{V} \bm{A}^T + c\bm{I}_p.
\end{equation*}
If $(\bm{V}, c) \in \ker(d\Phi_{(\bm{A},s)})$,
then

\begin{equation}\label{eq:ker_Phi_equation}
  \bm{A}\bm{V}^T + \bm{V} \bm{A}^T + c\bm{I}_p = \bm{0}.
\end{equation}
If $\bm{u} \in \mathbb{R}^p$ is any nonzero vector orthogonal to
all column vectors of $\bm{A}$, then $\bm{A}^T \bm{u} = \bm{0}$ and

\begin{equation*}
  0 = \bm{u}^T (\bm{A}\bm{V}^T + \bm{V} \bm{A}^T + c\bm{I}_p) \bm{u}
  = c \, \bm{u}^T \bm{u}.
\end{equation*}
Given that $\bm{u} \neq \bm{0}$, we must have $c = 0$
and \eqref{eq:ker_Phi_equation} reduces to

\begin{equation}\label{eq:ker_Phi_equation_reduced}
  \bm{A}\bm{V}^T + \bm{V} \bm{A}^T = \bm{0}.
\end{equation}
Apply $\bm{u}$ to the right to get:

\begin{equation*}
  \bm{A}\bm{V}^T\bm{u} = \bm{0}.
\end{equation*}
Given that $\bm{A}$ has full rank, 
rank nullity theorem forces 
\begin{equation*}
  \bm{V}^T\bm{u} = \bm{0}.
\end{equation*}
Therefore, we have proven that every vector $\bm{u}$
orthogonal to columns vectors of $\bm{A}$
will also be orthogonal to column vectors of $\bm{V}$.
This means $\bm{V}=\bm{A} \bm{\Omega}$ for some 
$q \times q$ matrix $\bm{\Omega}$.
Substitute this back to \eqref{eq:ker_Phi_equation_reduced}
gives 

\begin{equation*}
  \bm{A}(\bm{\Omega}^T + \bm{\Omega})\bm{A}^T = \bm{0}.
\end{equation*}
Applying $(\bm{A}^T \bm{A})^{-1} \bm{A}^T$ 
to the left and $\bm{A}(\bm{A}^T \bm{A})^{-1}$ to the right 
yields $\bm{\Omega}^T + \bm{\Omega} = \bm{0}$.
Therefore, $\bm{\Omega}$ is skew-symmetric.
Conversely, \eqref{eq:ker_Phi_equation_reduced} 
is satisfied by $\bm{V}=\bm{A} \bm{\Omega}$
with any skew-symmetric matrix $\bm{\Omega}$.
We have proven that $\ker(\mathrm{d}\Phi_{(\bm{A},s)})$
equals to the collection of all tangent vectors 
of the form $(\bm{A} \bm{\Omega}, 0) \in T_{(\bm{A},s)}\Theta$,
with $\bm{\Omega}$ ranged over all $q \times q$ skew symmetric matrices.

We now move back to prove 1 and 2 at the same time
by showing \eqref{eq:vertical_space}.
Clearly, $\bm{\Sigma}$ is well-defined 
and is injective.
The smoothness of $\bm{\Sigma}$
is guaranteed by relation $\bm{\Sigma} \circ \pi = \Phi$
(see for example \citealp{Lee_2013}, Theorem 4.29).
By the chain rule, we have 

\begin{equation} \label{eq:differential_Phi}
  \mathrm{d} \Phi_{(\bm{A},s)} = \mathrm{d} \bm{\Sigma}_{[\bm{A},s]} \circ \mathrm{d} \pi_{(\bm{A},s)}.
\end{equation}
It follows immediately that 
$\ker(\mathrm{d}\pi_{(\bm{A},s)}) \subseteq \ker(\mathrm{d}\Phi_{(\bm{A},s)})$.
The other direction will follow if we can prove that 
$\mathrm{d} \bm{\Sigma}_{[\bm{A},s]}$ is injective.

Suppose that $\bm{v} \in T_{[\bm{A},s]}M$ is 
such that $\mathrm{d} \bm{\Sigma}_{[\bm{A},s]} (\bm{v}) = 0$.
We prove that $\bm{v} = \bm{0}$.
Let $(\bm{V},c) \in T_{(\bm{A},s)}\Theta$ 
be the horizontal lift of $\bm{v}$ at $(\bm{A},s)$,
meaning that $(\bm{V},c)$ is the 
unique horizontal tangent vector at $(\bm{A},s)$
such that
 $\bm{v} = \mathrm{d} \pi_{(\bm{A},s)}((\bm{V},c))$.
\eqref{eq:differential_Phi} then implies that 
$(\bm{V},c) \in \ker(\mathrm{d} \Phi_{(\bm{A},s)})$
and the preceding discussion shows that 
$c = 0$ and $\bm{V} = \bm{A} \bm{\Omega}$ 
for some skew-symmetric $\bm{\Omega}$.
To proceed, 
we look at the 
curve $\gamma: (-\epsilon, \epsilon) \to \Theta$
defined by $\gamma(t) = (\bm{A}e^{t \bm{\Omega}}, s)$.
Clearly $\gamma(0)=(\bm{A}, s)$ and $\gamma'(0) = (\bm{A} \bm{\Omega}, 0)$.
Moreover, the matrix $e^{t \bm{\Omega}}$ is always orthogonal 
for any $t$ and any skew-symmetric $\bm{\Omega}$ 
so that $\pi(\gamma(t)) = [\bm{A},s]$ for all $t$.
We have 

\begin{equation*}
  \bm{v} =
  \mathrm{d} \pi_{(\bm{A},s)}((\bm{A} \bm{\Omega}, 0)) = 
  \mathrm{d} \pi_{(\bm{A},s)} (\gamma'(0)) = \frac{\mathrm{d}}{\mathrm{d} t}
   \bigg|_{t=0} \pi(\gamma(t)) = \bm{0} \in T_{[\bm{A},s]}M.
\end{equation*}
Therefore, $\mathrm{d} \bm{\Sigma}_{[\bm{A},s]}$ is injective and
the vertical tangent space is the same as $\ker(\mathrm{d}\Phi_{(\bm{A},s)})$.
Since $[\bm{A},s]$ is arbitrary, $\bm{\Sigma}$ is an immersion.

To compute the horizontal tangent space, any $(\bm{V}, c) \in V_{(\bm{A},s)}^{\perp}$ 
will satisfy 

\begin{equation}\label{eq:orthogonal_horizontal}
  0 = g_{(\bm{A},s)}^{\Theta} 
  \bracket{(\bm{V}, c), (\bm{A} \bm{\Omega}, 0)} = 
  \tr(\bm{A}^T \bm{V} \bm{\Omega})
\end{equation}
for all skew-symmetric $\bm{\Omega}$.
Decomposing $\bm{A}^T \bm{V} = \bm{S} + \bm{K}$, 
where $\bm{S}$ is the symmetric part and $\bm{K}$ is the skew-symmetric part,
we find out that 

\begin{equation*}
  0 = \tr(\bm{S} \bm{\Omega}) + \tr(\bm{K} \bm{\Omega}) = 0 + \tr(\bm{K} \bm{\Omega}) = \tr(\bm{K} \bm{\Omega}).
\end{equation*}
Setting $\bm{\Omega} = \bm{K}^T$ will yield that $\bm{K} = \bm{0}$.
Therefore, $\bm{A}^T \bm{V}$ is a symmetric matrix.
Conversely, any $\bm{V}$ such that $\bm{A}^T \bm{V}$ is symmetric will satisfy 
\eqref{eq:orthogonal_horizontal}. 
The proof is complete.
\end{proof}

We will need to compute the projection of a tangent 
vector to horizontal and verticle subspaces.

\begin{lemma}\label{thm:psi_S}
For any $q \times q$ symmetric positive definite matrix $\bm{S}$
and $q \times q$ matrix $\bm{\Omega}$,
define 

\begin{equation*}
  \psi_{\bm{S}}(\bm{\Omega}) = \bm{S} \bm{\Omega} + \bm{\Omega} \bm{S}.
\end{equation*}
Then, for each fixed $\bm{S}$, $\psi_{\bm{S}}$ is linear, 
maps skew-symmetric matrices to skew-symmetric matrices,
and is invertible. Furthermore, 
both $\psi_{\bm{S}}(\bm{\Omega})$
and $\psi_{\bm{S}}^{-1}(\bm{\Omega})$
depend smoothly on the tuple $(\bm{S}, \bm{\Omega})$.
\end{lemma}

\begin{proof}
It is clear from definition that $\psi_{\bm{S}}$ is linear and
$\psi_{\bm{S}}(\bm{\Omega})$ is skew-symmetric whenever $\bm{\Omega}$ is.
To show that $\psi_{\bm{S}}$ is invertible,
suppose that $\bm{\Omega}$ is such that 

\begin{equation*}
  \psi_{\bm{S}}(\bm{\Omega}) = \bm{0}.
\end{equation*}
Apply spectral theorem to get $\bm{Q} \in O(q)$ 
and a $q \times q$ diagonal matrix $\Lambda$ with positive entries 
$\lambda_1(\bm{S})$, $\ldots$, $\lambda_q(\bm{S})$
such that $\bm{S} = \bm{Q} \Lambda \bm{Q}^T$.
Substitution yields 
\begin{equation*}
  \Lambda \bm{Q}^T \bm{\Omega} \bm{Q} + \bm{Q}^T \bm{\Omega} \bm{Q} \Lambda = \bm{0}.
\end{equation*}
Looking at each entry, we arrive at 
\begin{equation*}
  (\lambda_i(\bm{S}) + \lambda_j(\bm{S})) 
(\bm{Q}^T \bm{\Omega} \bm{Q})_{ij} = 0
\end{equation*}
for $i,j = 1, \ldots, q$.
Given that $\bm{S}$ is positive definite, 
the eigenvalues of $\bm{S}$ are positive.
Thus, 
$\bm{Q}^T \bm{\Omega} \bm{Q} = \bm{0}$
and $\bm{\Omega} = \bm{0}$.
We have proved that $\psi_{\bm{S}}$ is invertible.
Given that $\psi_{\bm{S}}(\bm{\Omega})$ is a polynomial of the tuple $(\bm{S}, \bm{\Omega})$,
it depends smoothly on the tuple.
$\psi_{\bm{S}}^{-1}(\bm{\Omega})$ also depends smoothly on 
$(\bm{S}, \bm{\Omega})$ because the inverse is a smooth operation.
\end{proof}

\begin{proposition}\label{thm:projection}
Given $(\bm{A}, s) \in \Theta$ and $(\bm{V}, c) \in T_{(\bm{A}, s)}\Theta$,
orthogonal projection of $(\bm{V}, c)$ to $V_{(\bm{A}, s)}$ 
is

\begin{equation*}
  \bracket{\bm{A} \psi_{\bm{A}^T\bm{A}}^{-1}(\bm{A}^T \bm{V} - \bm{V}^T \bm{A}), 0}
\end{equation*}
and the same vector projected
to $H_{(\bm{A}, s)}$ is

\begin{equation}\label{eq:horizontal_projection}
  \bracket{\bm{V} - \bm{A} \psi_{\bm{A}^T\bm{A}}^{-1}(\bm{A}^T \bm{V} - \bm{V}^T \bm{A}), c}.
\end{equation}
Here $\psi$ is as in Lemma~\ref{thm:psi_S}.
\end{proposition}

\begin{proof}
It suffices to prove the formula for verticle projection 
as the horizontal one follows by a subtraction.
By Proposition~\ref{thm:horizontal_vertical},
the orthogonal projection of $(\bm{V}, c)$ to $V_{(\bm{A}, s)}$
equals to $(\bm{A} \bm{\Omega}, 0) \in V_{(\bm{A}, s)}$ for some $q \times q$ skew-symmetric matrix $\bm{\Omega}$
which minimizes $\|\bm{V}-	\bm{A} \bm{\Omega}\|_F^2$.
To compute $\bm{\Omega}$, let us compute the directional derivative of the quantity 
in the direction of a $q \times q$ skew-symmetric matrix $\bm{H}$.
For any $t \in \mathbb{R}$:

\begin{equation*}
  \begin{split}
  & \tab[0.4cm] \|\bm{V}-	\bm{A} (\bm{\Omega} + t\bm{H}) \|_F^2 
  = \tr\curly{\squared{\bm{V}-	\bm{A} (\bm{\Omega} + t\bm{H})}\squared{\bm{V}-	\bm{A} (\bm{\Omega} + t\bm{H})}^T}
  \\ & = \|\bm{V}\|_F^2 - 2 \tr(\bm{\Omega}^T\bm{A}^T\bm{V} ) - 2t \tr(\bm{H}^T\bm{A}^T\bm{V} )
  + \tr(\bm{\Omega}^T\bm{A}^T\bm{A}\bm{\Omega} )
  + t^2 \tr(\bm{H}^T\bm{A}^T\bm{A}\bm{H} )
  \end{split}
\end{equation*}
and 
\begin{equation*}
  \frac{\mathrm{d}}{\mathrm{d} t} \bigg|_{t=0}   \|\bm{V}-	\bm{A} (\bm{\Omega} + t\bm{H}) \|_F^2 
  = \tr \squared{\bm{H}^T (\bm{A}^T\bm{A}\bm{\Omega} - \bm{A}^T \bm{V})}.
\end{equation*}
Therefore, any solution to the equation setting all directional derivatives to $0$
yields that $\tr \squared{\bm{H}^T (\bm{A}^T\bm{A}\bm{\Omega} - \bm{A}^T \bm{V})} = 0$
for all skew-symmetric $\bm{H}$,
or equivalently that $\bm{A}^T\bm{A}\bm{\Omega} - \bm{A}^T \bm{V}$ is symmetric.
Therefore, $\Omega$ must satisfy 

\begin{equation*}
  \bm{A}^T\bm{A}\bm{\Omega} + \bm{\Omega} \bm{A}^T\bm{A} = \bm{A}^T \bm{V} - \bm{V}^T \bm{A}.
\end{equation*}
The proof is complete once we apply
Lemma~\ref{thm:psi_S}.
\end{proof}

In most textbooks, one can find the standard fact that 
any smooth vector field in $M$ uniquely horizontally lifts 
to a smooth vector field in $\Theta$.
However, not all textbooks mention about horizontal lifting curves
or vector fields along curves. We bookkeep such a proposition
for later convenience.

\begin{proposition}\label{thm:curve_lift}
Given a smooth curve $\gamma : (a,b) \to M$, 
a point $t_0 \in (a,b)$, and a representative $(\bm{A},s)$
of $\gamma(t_0)$, there exists a unique smooth curve $c: (a,b) \to \Theta$,
which we refer to as the horizontal lift of of $\gamma$ at $t_0$ and 
representative $(\bm{A},s)$,
such that $c(t_0) = (\bm{A},s)$ and, for all $t \in (a,b)$,
$\pi(c(t)) = \gamma(t)$, $c'(t) \in H_{c(t)}$, and $\mathrm{d} \pi_{c(t)}(c'(t)) = \gamma(t)$.
Furthermore,
given a smooth vector field $\bm{X}(t)$ along $\gamma$,
define, for each $t \in (a,b)$,
$\bm{Y}(t)$
to be the horizontal lift of $\bm{X}(t)$ at $c(t)$.
Then $\bm{Y}(t)$ is a smooth vector field along $c$.
\end{proposition}

\begin{proof}
The fact that $\gamma$ can be uniquely hozirontally lifted to a curve $c$
is immediate from Theorem 17.8 and 17.9 from \citealp{Michor_2008}.
Given a smooth vector field $\bm{X}(t)$ along $\gamma$
and any $t \in (a,b)$, we can choose a smooth local frame 
around $\gamma(t)$ so that the coefficients of $\bm{X}(t)$ 
under this frame is smooth.
Given that smooth local frames in $M$ lifts horizontally 
to smooth local frames in $\Theta$, 
the coefficients of $\bm{Y}(t)$ is defined as in the proposition
will be smooth under the lifted frame, meaning that 
$\bm{Y}(t)$ is smooth around $t$.
Since $t$ is arbitrary, we have proved that $\bm{Y}(t)$
is a smooth vector field along $c$.

Given that the notations and language are slightly different in \citealp{Michor_2008},
we provide a bit more explanation for reader's convenience. 
We apply the two theorems for the fiber bundle $(E, p, M, S)$ with $E = \Theta$, $p = \pi$, $S = O(q)$,
and $M$ is just our $M$. 
The map $\Phi$ mentioned in their text (called the Ehresmann connection),
merely assigns each $\theta \in \Theta$ 
the linear projection map $\Phi_{\theta}:T_{\theta}\Theta \to T_{\theta}\Theta$
onto the subspace $V_{\theta}$.
Given the smooth curve $\gamma$ in the statement of our proposition,
Theorem 17.8 implemented with $x = \gamma(t_0)$ and $u_x = (\bm{A},s)$
produces a map $\text{Pt}_{\gamma}$.
We set $c(t) = \text{Pt}_{\gamma}((\bm{A},s),t)$, which is a smooth curve.
Property (1) of Theorem 17.8 shows that $c(t_0) = (\bm{A},s)$
and $\pi(c(t)) = \gamma(t)$ for all $t$ such that $c$ is well-defined.
Property (2) of Theorem 17.8 shows that $c'(t) \in H_{c(t)}$ 
for all $t$ such that $c$ is well-defined.
Given that $S = O(q)$ is compact, Theorem 17.9
shows that $c$ is defined on the entire $(a,b)$.
An application of chain rule from $\pi \circ c = \gamma$
shows that $\mathrm{d} \pi_{c(t)}(c'(t)) = \gamma(t)$
for all $t \in (a,b)$.
Finally, as pointed out in the first proof of Theorem 17.8,
$\text{Pt}_{\gamma}$ is unique, or equivalently $c$ is unique.
\end{proof}

\subsection{Geodesics, the Exponential Map, and the Logarithmic Map}
Here is a classification result about all geodesics in $M$.

\begin{proposition}\label{thm:geodesic_classification}
Given $[\bm{A},s] \in M$ and $\bm{v} \in T_{[\bm{A},s]}M$, 
all geodesics in $M$ starting at $[\bm{A},s]$ with initial velocity 
$\bm{v}$ takes the format of curves 
$\gamma: I \to M$, where $I$ is an interval containing $0$ in its interior,
defined by 
\begin{equation}\label{eq:geodesic_gamma}
  \gamma(t) = [\bm{A} + t\bm{V}, s + tc],
\end{equation}
where $(\bm{V},c) \in H_{(\bm{A},s)}$ is the horizontal lift 
of $\bm{v}$ at $ (\bm{A},s)$.
\end{proposition}

\begin{proof}
First of all, the curve $\gamma$ given by \eqref{eq:geodesic_gamma} is well defined. 
If $(\bm{A}\bm{Q},s)$, where $\bm{Q} \in O(q)$, is another representative of $[\bm{A},s]$,
let $(\bm{V}',c')$ be the horizontal lift of $\bm{v}$ to $(\bm{A}\bm{Q},s)$.
With the function $R_{\bm{Q}}: \Theta \to \Theta$ defined by $R_{\bm{Q}}((\bm{A},s)) = (\bm{A}\bm{Q},s)$,
we have 

\begin{equation*}
  \bm{v} = \mathrm{d}\pi_{(\bm{A},s)} (\bm{V},c) = \mathrm{d}(\pi \circ R_{\bm{Q}})_{(\bm{A},s)} (\bm{V},c)
  = \mathrm{d} \pi_{(\bm{A}\bm{Q},s)} (dR_{\bm{Q}})_{(\bm{A},s)}(\bm{V},c)
  = \mathrm{d} \pi_{(\bm{A}\bm{Q},s)} (\bm{V}\bm{Q},c).
\end{equation*}
Thus, $\bm{V}' = \bm{V}\bm{Q}$ and $c'=c$.
Substitute this into the definition of $\gamma$ 
will yield the exact same curve.

Given that $\gamma$ is well-defined, 
it must be that $(\bm{A} + t\bm{V}, s + tc) \in \Theta$ for all $t \in I$.
The curve $c:I \to \Theta$
with 

\begin{equation}\label{eq:geodesic_c}
  c(t) = (\bm{A} + t\bm{V}, s + tc)
\end{equation}
for all $t \in I$
is then well-defined. Being a straight line 
in Euclidean space, $c$ is a geodesic.
Moreover, it has horizontal initial velocity and satisfies 
$\pi \circ c = \gamma$. 
By Corollary 26.12 from \citealp{Michor_2008},
$\gamma$ must be a geodesic.
Conversely, suppose that $\gamma: I \to M$ is a geodesic in $M$
starting at
$[\bm{A},s]$ with initial velocity 
$\bm{v}$. 
Let $c$ be the horizontal lift of $\gamma$  at $0$ and representative 
$ (\bm{A},s)$
given by Proposition~\ref{thm:curve_lift}.
Then $c$ starts at $ (\bm{A},s)$ and has its initial velocity 
$(\bm{V},c)$ being the horizontal lift of $\bm{v}$ at $ (\bm{A},s)$.
Lemma 26.11 in \citealp{Michor_2008}
shows that $c$ is a geodesic in $\Theta$.
Given that straight lines are the only geodesics in Euclidean spaces,
$c$ must be defined by \eqref{eq:geodesic_c} for all $t \in I$
and $\gamma = \pi \circ c$ must take the format \eqref{eq:geodesic_gamma}.
\end{proof}

We now move to compute the exponential map 
and the logarithmic map.
The following two lemmas are crucial 
for computing the orthogonal matrix based on given points in $M$.

\begin{lemma}\label{thm:q_classification}
If $\bm{L}$ is an invertible $q \times q$ matrix
and $\bm{Q} \in O(q)$ are such that $\bm{L}\bm{Q}$ is symmetric,
then there exists matrices $\bm{R}\in O(q)$, $\bm{P} \in O(q)$, 
a $q \times q$ diagonal matrix $\bm{D}$ whose diagonal entries are $1$ or $-1$,
and a $q \times q$ diagonal matrix $\bm{S}$ with diagonal entries 
$\sigma_1(\bm{L}) \geqslant \cdots \geqslant \sigma_q(\bm{L})>0$
such that $\bm{L} = \bm{R}\bm{S}\bm{P}^T$
and $\bm{Q} = \bm{P}\bm{D}\bm{R}^T$.

If we add furthermore the assumption that 
$\bm{L}$ is symmetric positive definite,
then we can achieve the results mentioned above with $\bm{P} = \bm{R}$.
In particular, this implies 
$\bm{Q} = \bm{Q}^T$ and $\|\bm{Q}-\bm{I}_q\| = \|\bm{D}-\bm{I}_q\|$
can take only $0$ or $2$ as values.
\end{lemma}

\begin{proof}
Let us start by applying the singular value decomposition 
of $\bm{L}$ to get matrices $\bm{R}'\in O(q)$, $\bm{P}' \in O(q)$
and a diagonal matrix $\bm{S}$ with diagonal entries 
$\sigma_1(\bm{L}) \geqslant \cdots \geqslant \sigma_q(\bm{L})>0$ 
such that $\bm{L} = \bm{R}'\bm{S}(\bm{P}')^T$.
Note that all singular values are positive because 
$\bm{L}$ is invertible.
Substituting the given condition that $\bm{L}\bm{Q}$ is symmetric will yield that
$\bm{S}\bm{J} = \bm{J}^T\bm{S}$, where $\bm{J} := (\bm{P}')^T \bm{Q} \bm{R}'$.

Clearly $\bm{J} \in O(q)$.
The equation 

\begin{equation*}
  \bm{S}^2 = \bm{S}\bm{J}\bm{J}^T\bm{S} = (\bm{S}\bm{J})(\bm{S}\bm{J})^T 
  = (\bm{J}^T\bm{S})(\bm{J}^T\bm{S})^T = \bm{J}^T\bm{S}^2\bm{J}
\end{equation*}
shows that 
\begin{equation*}
  \bm{J}\bm{S}^2=\bm{S}^2\bm{J}.
\end{equation*}
Looking at each entry of both side,
we end up with 
\begin{equation*}
  \bm{J}_{ij}(\sigma_j(\bm{L})^2-\sigma_i(\bm{L})^2) = 0
\end{equation*}
for $i,j = 1, \ldots q$.
Given that all singular values are positive,
we in fact have
\begin{equation}\label{eq:j_s_equation}
  \bm{J}_{ij}(\sigma_j(\bm{L})-\sigma_i(\bm{L})) = 0
\end{equation}
for all $i$ and all $j$ so that $\bm{J}\bm{S}=\bm{S}\bm{J}=\bm{J}^T\bm{S}$.
Inverting $\bm{S}$ yields that $\bm{J}$ is actually symmetric.
Being a symmetric orthogonal matrix, $\bm{J}^2 = \bm{I}_q$ and the eigenvalues of $\bm{J}$ 
must be $1$ or $-1$.
Another look at \eqref{eq:j_s_equation} shows that $\bm{J}$ is 
a block-diagonal matrix, diagonal if 
all singular values are distinct.
Furthermore, each block of $\bm{J}$ must still be orthogonal and symmetric.

We now diagonalize $\bm{J}$ based on the above information.
Suppose that $\bm{S} = \oplus_{i=1}^k \bm{s}_i\bm{I}_{m_i}$,
where $\bm{s}_1 > \ldots > \bm{s}_k$ are the unique diagonal entries of 
$\bm{S}$ and each repeats for $m_i$ times.
Thus, $\bm{J}$ has $k$ blocks.
Write the $i$th block by $\bm{J}_i$ so that
$\bm{J} = \oplus_{i=1}^k \bm{J}_i$.
The spectral theorem then allows us to orthogonally diagonalize 
$\bm{J}_i = \bm{K}_i\bm{D}_i\bm{K}_i^T$,
where 
$\bm{K}_i \in O(m_i)$ and
$\bm{D}_i$ is an $m_i \times m_i$ diagonal matrix with 
all diagonal entries being $1$ or $-1$. 
Put $\bm{K} = \oplus_{i=1}^k \bm{K}_i$
and $\bm{D} = \oplus_{i=1}^k \bm{D}_i$.
Then $\bm{K} \in O(q)$, $\bm{D}$ is a $q \times q$ diagonal matrix 
with $1$ or $-1$ diagonal entry, and $\bm{J} = \bm{K}\bm{D}\bm{K}^T$.
This involved construction leads us to the outcome that 
$\bm{K}$ in fact commutes with $\bm{S}$:

\begin{equation*}
  \bm{K}\bm{S} = \bigoplus_{i=1}^k \bm{s}_i\bm{K}_i = \bm{S}\bm{K}.
\end{equation*}
Put $\bm{P} = \bm{P}'\bm{K}$,
$\bm{R} = \bm{R}'\bm{K}$.
We have 
\begin{equation*}
  \bm{Q} = \bm{P}' \bm{J} (\bm{R}')^T = \bm{P}' \bm{K}\bm{D}\bm{K}^T (\bm{R}')^T
  = \bm{P} \bm{D} \bm{R}^T
\end{equation*}
and
\begin{equation*}
  \bm{R}\bm{S}\bm{P}^T = 
  \bm{R}'\bm{K} \bm{S}\bm{K}^T(\bm{P}')^T
  = \bm{R}'\bm{S} \bm{K}\bm{K}^T(\bm{P}')^T
  = \bm{R}'\bm{S} (\bm{P}')^T =\bm{L}.
\end{equation*}
This completes the proof.

Finally, let us just note that, in case 
$\bm{L}$ is symmetric positive definite,
the spectral theorem allows us to take $\bm{R}'=\bm{P}'$ from the start 
and it will yield in the end 
$\bm{L} = \bm{P}\bm{S}\bm{P}^T$
and $\bm{Q} = \bm{P}\bm{D}\bm{P}^T$.
It then follows easily that $\bm{Q}^2 = \bm{I}_q$
and $\bm{Q}^T = \bm{Q}$. If $\bm{D} = \bm{I}_q$,
then $\|\bm{Q}-\bm{I}_q\| = \|\bm{D}-\bm{I}_q\|=0$.
If any entry of $\bm{D}$ is $-1$, $\|\bm{D}-\bm{I}_q\| = 2$.
\end{proof}

\begin{lemma}\label{thm:unique_P_R}
If $\bm{L}$ is an invertible $q \times q$ matrix
with singular value decomposition $\bm{L} = \bm{R}\bm{S}\bm{P}^T$,
then 

\begin{equation*}
  \bm{P}\bm{R}^T = (\bm{L}^T \bm{L})^{-\frac{1}{2}}\bm{L}^T.
\end{equation*}
\end{lemma}

\begin{proof}
Given that

\begin{equation*}
  \bm{L}^T\bm{L} = \bm{P}\bm{S}^2\bm{P}^T,
\end{equation*}
we have
\begin{equation*}
  (\bm{L}^T \bm{L})^{\frac{1}{2}} = \bm{P}\bm{S}\bm{P}^T.
\end{equation*}
Therefore:
\begin{equation*}
  \bm{L} = \bm{R}\bm{P}^T\bm{P}\bm{S}\bm{P}^T = \bm{R}\bm{P}^T (\bm{L}^T \bm{L})^{\frac{1}{2}}.
\end{equation*}
The result follows from inverting $(\bm{L}^T \bm{L})^{\frac{1}{2}}$ and take transpose.
\end{proof}

We now state and prove all the results about 
the exponential and the logarithmic map 
in one single proposition.

\begin{proposition}\label{thm:exponential_map}
The exponential map 
is computed as 

\begin{equation}\label{eq:exponential_map}
  \text{Exp}_{[\bm{A},s]}(\bm{v}) = [\bm{A} + \bm{V}, s + c],
\end{equation}
where $[\bm{A},s] \in M$, $\bm{v} \in T_{[\bm{A},s]}M$,
and $(\bm{V},c)$ is the horizontal lift of $\bm{v}$ at $(\bm{A},s)$.

Let us use $\mathcal{E}_{[\bm{A},s]}$
to denote the open ball in $T_{[\bm{A},s]}M \cong \mathbb{R}^{d}$
centred at $\bm{0}$ of radius $\min(\sigma_{\text{min}}(\bm{A}),s)$.
Given that \eqref{eq:exponential_map} is well-defined 
at least
over all $\bm{v} \in \mathcal{E}_{[\bm{A},s]}$,
we from now on treat $\mathcal{E}_{[\bm{A},s]}$
as the domain of the exponential map.
That is, the function $\text{Exp}_{[\bm{A},s]}: \mathcal{E}_{[\bm{A},s]} \to U_{[\bm{A},s]}$,
where $U_{[\bm{A},s]} := \text{Exp}_{[\bm{A},s]}(\mathcal{E}_{[\bm{A},s]})$ is an 
open subset of $M$ referred to as a geodesic ball,
is defined by \eqref{eq:exponential_map}.
Note the following facts:

\begin{enumerate}
  \item For every $\bm{v} \in \mathcal{E}_{[\bm{A},s]}$,
  $\bm{A}^T(\bm{A}+\bm{V})$ is symmetric positive definite.
  \item $\text{Exp}_{[\bm{A},s]}$ is a diffeomorphism.
  \item If $[\bm{B},t] \in U_{[\bm{A},s]}$, then for every representative $(\bm{A},s)$ of $[\bm{A},s]$
  and every representative $(\bm{B},t)$ of $[\bm{B},t]$, the matrix $\bm{A}^T \bm{B}$ will 
  always be invertible. Given representative $(\bm{A},s)$ of $[\bm{A},s]$,
  \textit{some} representative $(\bm{B},t)$ of $[\bm{B},t]$ will furthermore 
  make $\bm{A}^T \bm{B}$ symmetric positive definite.
  \item The logarithmic map $\text{Exp}^{-1}_{[\bm{A},s]}:U_{[\bm{A},s]} \to \mathcal{E}_{[\bm{A},s]}$
  is computed as follows:
  for every $[\bm{B},t] \in U_{[\bm{A},s]}$,
  \begin{equation}\label{eq:logarithmic_def}
    \text{Exp}^{-1}_{[\bm{A},s]}([\bm{B},t]) 
    = \mathrm{d} \pi_{(\bm{A},s)}((\bm{B}\bm{Q}(\bm{A},\bm{B})-\bm{A}, t-s)),
  \end{equation}
  where 

  \begin{equation}\label{eq:Q_A_B_def}
    \bm{Q}(\bm{A},\bm{B}) 
    = (\bm{B}^T \bm{A} \bm{A}^T \bm{B})^{- \frac{1}{2}} \bm{B}^T \bm{A} \in O(q).
  \end{equation}
  \item The geodesic ball $U_{[\bm{A},s]}$ equals to the set of all $[\bm{B},t] \in M$ such that 
  $\bm{A}^T\bm{B}$ is invertible and
  \begin{equation}\label{eq:U_A_s_def}
    \|\bm{B}\bm{Q}(\bm{A},\bm{B})-\bm{A}\|_F^2 + (t-s)^2
    < \min\bracket{\sigma_{\min}(\bm{A}),s}^2.
  \end{equation}
  \item If $[\bm{B},t] \in U_{[\bm{A},s]}$, then the Riemannian distance between 
  $[\bm{A},s]$ and $[\bm{B},t]$ equals to 

  \begin{equation}\label{eq:Riemannian_distance}
    \abs{\text{Exp}^{-1}_{[\bm{A},s]}([\bm{B},t])} = \sqrt{\|\bm{B}\bm{Q}(\bm{A},\bm{B})-\bm{A}\|_F^2 + (t-s)^2},
  \end{equation}
  which is realized by one unique geodesic $\gamma: [0,1] \to M$ defined by, for all $u \in [0,1]$:
  \begin{equation}\label{eq:dist_minimizing_curve}
    \gamma(u) = \text{Exp}_{[\bm{A},s]}(u \text{Exp}^{-1}_{[\bm{A},s]}([\bm{B},t])).
  \end{equation}
\end{enumerate}
\end{proposition}

\begin{proof}
Before we start the proof, let us first check that all the 
statements are indeed well-defined.
If $\bm{R} \in O(q)$ and $\bm{P} \in O(q)$,
then 
\begin{equation}\label{eq:Q_A_B_RP}
  \bm{Q}(\bm{A}\bm{P},\bm{B}\bm{R}) = \bm{R}^T\bm{Q}(\bm{A},\bm{B})\bm{P}.
\end{equation}
Therefore,
the same computation as in the proof of Proposition~\ref{thm:geodesic_classification} will show that 
\eqref{eq:logarithmic_def} is indeed well-defined.
Also, the contents of the inequality \eqref{eq:U_A_s_def} 
and the expression \eqref{eq:Riemannian_distance}
are independent of representatives and the description of 
the set $U_{[\bm{A},s]}$ makes sense.

Proposition~\ref{thm:geodesic_classification} already shows 
that the computation of the exponential map is valid as long as 
$(\bm{A} + \bm{V}, s + c) \in \Theta$.
Therefore, the exponential map is certainly well-defined 
over all $\bm{v}$ with small enough norm.
To have a quantified upper bound of the norm,
we use the classical result of
Weyl's perturbation inequality 
(see for example Chapter 7.3 Problem 16 equation (7.3.15) of \citealp{Horn_Johnson_2013})
to get

\begin{equation*}
  \sigma_{\text{min}}(\bm{A} + \bm{V}) \geqslant \sigma_{\text{min}}(\bm{A}) - \|\bm{V}\|
  \geqslant \sigma_{\text{min}}(\bm{A}) - \|\bm{V}\|_F \geqslant \sigma_{\text{min}}(\bm{A}) - \abs{\bm{v}} > 0,
\end{equation*}
whenever $\abs{\bm{v}} < \min(\sigma_{\text{min}}(\bm{A}),s)$
or equivalently $\bm{v} \in \mathcal{E}_{[\bm{A},s]}$.
We have made use of the two facts that 
$\|\bm{V}\| \leqslant \|\bm{V}\|_F$
for any matrix $\bm{V}$
and that $\abs{\bm{v}}=\sqrt{\|\bm{V}\|_F^2+c^2}$
by isometry.
Thus, for all $\bm{v} \in \mathcal{E}_{[\bm{A},s]}$, $\rank(\bm{A} + \bm{V}) = q$
and the exponential map is well-defined in $\mathcal{E}_{[\bm{A},s]}$.

We now start the proof of the five facts 
of the function $\text{Exp}_{[\bm{A},s]}: \mathcal{E}_{[\bm{A},s]} \to U_{[\bm{A},s]}$.

The first fact can be quickly settled.
Let $\bm{v} \in \mathcal{E}_{[\bm{A},s]}$
with horizontal lift $(\bm{V},c)$ to $(\bm{A},s)$ be given.
As proved in Proposition~\ref{thm:horizontal_vertical},
$\bm{V}^T\bm{A} = \bm{A}^T\bm{V}$ and thus
$\bm{A}^T(\bm{A}+\bm{V})$ is symmetric.
If $\bm{x} \in \mathbb{R}^q$ is any unit vector, then 

\begin{equation*}
  \bm{x}^T \bm{A}^T(\bm{A}+\bm{V}) \bm{x}
  = \abs{\bm{A}\bm{x}}^2 + \bracket{\bm{A}\bm{x}}^T \bm{V}\bm{x}
  \geqslant \abs{\bm{A}\bm{x}}(\abs{\bm{A}\bm{x}}-\abs{\bm{V}\bm{x}})
  \geqslant \sigma_{\min}(\bm{A})(\sigma_{\min}(\bm{A})-	\|\bm{V}	\|_F) > 0
\end{equation*}
because $\sigma_{\min}(\bm{A}) >	\abs{\bm{v}} \geqslant \|\bm{V}\|_F$.

The second fact has a longer proof.
Being a well-defined exponential map,
$\text{Exp}_{[\bm{A},s]}$ is automatically a smooth map.
To prove that it is actually a diffeomorphism 
in $\mathcal{E}_{[\bm{A},s]}$,
we must first show that it is injective in 
$\mathcal{E}_{[\bm{A},s]}$ so that 
the inverse function, the logarithmic map $\text{Exp}^{-1}_{[\bm{A},s]}$,
exists. Once this is done, 
we prove that $U_{[\bm{A},s]}$ is open and 
that $\text{Exp}^{-1}_{[\bm{A},s]}$ 
is smooth in one shot 
by proving that its differential
$\mathrm{d} \bracket{\text{Exp}_{[\bm{A},s]}}_{\bm{v}}$
is an isomorphism for every given $\bm{v} \in \mathcal{E}_{[\bm{A},s]}$.

We thus move to show that $\text{Exp}_{[\bm{A},s]}$
is injective in $\mathcal{E}_{[\bm{A},s]}$.
Suppose that $\bm{v}_1$ and $\bm{v}_2$
are two vectors in $\mathcal{E}_{[\bm{A},s]}$
who lifts to $(\bm{V}_1,c_1)$ and $(\bm{V}_2,c_2)$
respectively 
at $(\bm{A},s)$
such that $\text{Exp}_{[\bm{A},s]}(\bm{v}_1) = \text{Exp}_{[\bm{A},s]}(\bm{v}_2)$.
We prove that $\bm{v}_1=\bm{v}_2$.
\eqref{eq:exponential_map} immediately
shows that $c_1=c_2$ and 

\begin{equation*}
  \bm{A} + \bm{V}_1 = (\bm{A} + \bm{V}_2)\bm{Q}
\end{equation*}
for some $\bm{Q} \in O(q)$.
Lemma~\ref{thm:q_classification}
with the symmetric positive definite matrix $\bm{L} = \bm{A}^T (\bm{A} + \bm{V}_2)$
shows that $\|\bm{Q} - \bm{I}_q\|$ can only equal to $0$ or $2$.
Given that

\begin{equation*}
  \sigma_{\min}(\bm{A})\|\bm{Q} - \bm{I}_q\| \leqslant 
  \|\bm{A}\bm{Q} - \bm{A}\| = \|\bm{V}_1 - \bm{V}_2\bm{Q}\|
  \leqslant \|\bm{V}_1\| + \| \bm{V}_2\| < 2\sigma_{\min}(\bm{A}),
\end{equation*}
we have proved that $\|\bm{Q} - \bm{I}_q\| < 2$ and $\bm{Q} = \bm{I}_q$.
Therefore, $\bm{V}_1 =\bm{V}_2$ and $\bm{v}_1=\bm{v}_2$.
The exponential map is indeed injective in $\mathcal{E}_{[\bm{A},s]}$.

Next, let $\bm{v} \in \mathcal{E}_{[\bm{A},s]}$
be given. We prove that $\mathrm{d} \bracket{\text{Exp}_{[\bm{A},s]}}_{\bm{v}}$
is an isomorphism because $\ker \bracket{\mathrm{d} \bracket{\text{Exp}_{[\bm{A},s]}}_{\bm{v}}}$
consists only of the zero vector.
Suppose that $\bm{v}$ lifts to $(\bm{V},c)$ at $(\bm{A},s)$.
To make the task notationally feasible, we 
implicitly canonically identify the Euclidean spaces together.
Let us write $l: T_{[\bm{A},s]}M \to H_{(\bm{A},s)}$ 
for the horizontal lift map: $l(\bm{v}) = (\bm{V},c)$,
which is a linear isometry.
Define a map $F: l(\mathcal{E}_{[\bm{A},s]}) \to M$
by setting, for 
any $(\bm{U},b) \in l(\mathcal{E}_{[\bm{A},s]}) \subseteq H_{(\bm{A},s)}$,
that $F((\bm{U},b)) = \pi((\bm{A}+\bm{U}, s + b))$.
In this notation, $\text{Exp}_{[\bm{A},s]} = F \circ l$
and the chain rule shows that 

\begin{equation*}
  \mathrm{d} \bracket{\text{Exp}_{[\bm{A},s]}}_{\bm{v}} = 
  \mathrm{d} F_{(\bm{V},c)} \circ \mathrm{d} l_{\bm{v}}.
\end{equation*}
Therefore, the only thing left to show is that 
$\ker \bracket{\mathrm{d} F_{(\bm{V},c)}}$ contains only the zero vector.
For any $(\bm{U},b) \in l(\mathcal{E}_{[\bm{A},s]})$,
we always have 

\begin{equation*}
  \mathrm{d} F_{(\bm{V},c)} ((\bm{U},b))
  = \frac{d}{dt}\bigg|_{t=0} \pi((\bm{A}+\bm{V}+t\bm{U}, s+c+tb))
  = \mathrm{d} \pi_{(\bm{A}+\bm{V}, s+c)}((\bm{U},b)).
\end{equation*}
Therefore, if $(\bm{U},b) \in \ker \bracket{\mathrm{d} F_{(\bm{V},c)}}$,
then $(\bm{U},b) \in V_{(\bm{A}+\bm{V}, s+c)}$
so that $b = 0$ and $\bm{U} = (\bm{A}+\bm{V})\bm{\Omega}$
for some $q \times q$ skew-symmetric matrix $\bm{\Omega}$.
But $(\bm{U},b) \in H_{(\bm{A}, s)}$ as well.
Thus, $\bm{A}^T(\bm{A}+\bm{V})\bm{\Omega} = \bm{\Omega}^T (\bm{A}+\bm{V})^T \bm{A}$.
Given that $\bm{A}^T(\bm{A}+\bm{V})$ is also symmetric,
we arrive at 

\begin{equation*}
  \bm{A}^T(\bm{A}+\bm{V})\bm{\Omega} + \bm{\Omega} \bm{A}^T(\bm{A}+\bm{V}) =\bm{0}.
\end{equation*}
Lemma~\ref{thm:psi_S} with $\bm{S} = \bm{A}^T(\bm{A}+\bm{V})$
forces $\bm{\Omega} = \bm{0}$.
We have proven that $\text{Exp}_{[\bm{A},s]}$
is a diffeomorphism in $\mathcal{E}_{[\bm{A},s]}$
and that $U_{[\bm{A},s]}$ is an open subset of $M$.

Now that the logarithmic map
$\text{Exp}^{-1}_{[\bm{A},s]}:U_{[\bm{A},s]} \to \mathcal{E}_{[\bm{A},s]}$
is well-defined and is smooth,
we will compute it and 
show the next three facts 
in one shot.
Let $[\bm{B},t] \in U_{[\bm{A},s]}$ be given.
By \eqref{eq:exponential_map},
$[\bm{B},t] = [\bm{A} + \bm{V},s+c]$,
where $(\bm{V}, c)$ is the horizontal lift of 
$\bm{v} = \text{Exp}^{-1}_{[\bm{A},s]}([\bm{B},t])$ to $(\bm{A}, s)$.
It immediately follows that $c = t - s$.
The first fact shows that, with $\bm{B}_0 := \bm{A} + \bm{V}$,
$\bm{A}^T\bm{B}_0$
is symmetric positive definite.
Given that the value 
of $\abs{\det(\bm{A}^T\bm{B})}$
is independent of representatives,
$\bm{A}^T\bm{B}$ is always invertible 
for any representative.
The third fact is therefore established.
Given that $\bm{Q}(\bm{A},\bm{B}_0) = \bm{I}_q$,
we immediately see that \eqref{eq:logarithmic_def} 
holds with $\bm{B}_0$ 
in place of $\bm{B}$.
However, \eqref{eq:logarithmic_def} is independent 
of representatives so it is valid for every representative $\bm{B}$.
Finally, 
let us write $C$ for the set of all points $[\bm{B},t] \in M$
such that $\bm{A}^T \bm{B}$ is invertible and
satisfies \eqref{eq:U_A_s_def}.
By using $\bm{B}_0$ 
in place of $\bm{B}$ in \eqref{eq:U_A_s_def},
it is clear that
$U_{[\bm{A},s]} \subseteq C$.
Conversely, let $[\bm{B},t] \in C$ be given.
Put $\bm{B}_0:=\bm{B}\bm{Q}(\bm{A}, \bm{B})$.
Then $\bm{A}^T\bm{B}_0$ is again symmetric 
positive definite.
In particular, $\bm{A}^T(\bm{B}_0-\bm{A})$
is symmetric and $(\bm{V}, c) := (\bm{B}_0-\bm{A}, t-s)$
is a horizontal vector at $(\bm{A},s)$.
If $\bm{v} := \mathrm{d} \pi_{(\bm{A},s)}((\bm{V}, c))$,
then \eqref{eq:U_A_s_def} shows that $\bm{v} \in \mathcal{E}_{[\bm{A},s]}$
and $[\bm{B},t] \in U_{[\bm{A},s]}$.
We conclude that $U_{[\bm{A},s]} = C$.

Finally, let us calculate the Riemannian distance 
between $[\bm{B},t] \in U_{[\bm{A},s]}$ and $[\bm{A},s]$.
Proposition 6.11 in \cite{Lee_2018} shows that
the Riemannian distance between points in a geodesic ball is indeed attained 
by exactly one geodesic.
It suffices to find a geodesic minimizing the Riemannian distance 
among all geodesics.
By Proposition~\ref{thm:geodesic_classification},
all geodesics connecting $[\bm{A},s]$ to $[\bm{B},t]$
takes the form $\gamma(u)=[\bm{A}+u\bm{V},s+uc]$,
where $(\bm{V},c)$ is the horizontal lift of some 
vector $\bm{v} \in T_{[\bm{A},s]}M$.
It is clear that $c = t-s$ and $\bm{V} = \bm{B}\bm{Q}-\bm{A}$
for some $\bm{Q} \in O(q)$.
Given that the endpoints are fixed,
the time is fixed,
and $\gamma$ has constant speed $\abs{\bm{v}}$,
the choice of $\bm{v}$ with minimal speed 
will give the smallest arclength and hence the smallest Riemannian distance.
To carry out this minimization,
Lemma~\ref{thm:q_classification} shows that
there exists $\bm{R} \in O(q)$, $\bm{P} \in O(q)$
such that
$\bm{A}^T\bm{B} = \bm{R}\bm{S}\bm{P}^T$
and $\bm{Q} = \bm{P}\bm{D}\bm{R}^T$.
Let us compute the norm of $\bm{v}$:
\begin{align*}
  & \tab[0.4cm] g_{[\bm{A},s]}(v, v)
  = g_{(\bm{A},s)}^{\Theta}\bracket{(\bm{B}\bm{Q}-\bm{A}, t-s), (\bm{B}\bm{Q}-\bm{A}, t-s)} \\
  &= \tr\bracket{(\bm{B}\bm{Q}-\bm{A})(\bm{B}\bm{Q}-\bm{A})^T} + (t-s)^2 
  = \tr(\bm{B}\bm{B}^T) + \tr(\bm{A}\bm{A}^T) - 2\tr(\bm{S}\bm{D}) + (t-s)^2,
\end{align*}
where the last equality substituted
$\bm{A}^T\bm{B}\bm{Q}' = \bm{R}\bm{S}\bm{P}^T\bm{P}\bm{D}\bm{R}^T = \bm{R}\bm{S}\bm{D}\bm{R}^T$.
The fact that
the diagonal entries of $\bm{S}$ are strictly positive
show that
$g_{[\bm{A},s]}(v, v)$ is minimized precisely when $\bm{D} = \bm{I}_q$.
By Lemma~\ref{thm:unique_P_R},
it must be that $\bm{Q} = \bm{P}\bm{R}^T = \bm{Q}(\bm{A}, \bm{B})$
and the last fact holds.
The proof is complete.
\end{proof}

\begin{remark}
In the above proof,
it might seem that
the proof of
\eqref{eq:logarithmic_def}
is trivialized while
the amazing matrix
\eqref{eq:Q_A_B_def} 
seems to come from nowhere.
This is not at all the case.
Originally,
the matrix $\bm{Q}(\bm{A},\bm{B})$ comes from 
considering the cases for a generic representative whose matrix 
$\bm{A}^T\bm{B}$
is only invertible, but not necessarily symmetric or positive definite.
Given
$[\bm{B},t] \in U_{[\bm{A},s]}$,
\eqref{eq:exponential_map} 
shows that 
\begin{equation*}
    \text{Exp}^{-1}_{[\bm{A},s]}([\bm{B},t]) 
    = \mathrm{d} \pi_{(\bm{A},s)}((\bm{B}\bm{Q}-\bm{A}, t-s))
  \end{equation*}
for some $\bm{Q}\in O(q)$.
Given that $\bm{A}^T(\bm{B}\bm{Q}-\bm{A})$ is symmetric,
it follows that $\bm{A}^T\bm{B}\bm{Q}$
is symmetric and Lemma~\ref{thm:q_classification}
shows that 
there exists $\bm{R} \in O(q)$, $\bm{P} \in O(q)$
such that
$\bm{A}^T\bm{B} = \bm{R}\bm{S}\bm{P}^T$
and $\bm{Q} = \bm{P}\bm{D}\bm{R}^T$.
Thus, by continuity of the logarithmic map,
if $[\bm{B},t] \to [\bm{A},s]$ and for some 
representative 
$(\bm{B},t) \to (\bm{A},s)$,
then
it must be the case that $\bm{B}\bm{Q}-\bm{A} \to \bm{0}$
or equivalently $\bm{Q} \to \bm{I}_q$.
Therefore, heurestically, the only reasonable choice of 
$\bm{D}$ is $\bm{I}_q$.
So $\bm{Q} = \bm{P}\bm{R}^T$.
This value is actually uniquely determined 
by $\bm{A}$ and $\bm{B}$ and gives exactly 
$\bm{Q} = \bm{Q}(\bm{A},\bm{B})$, as shown in
Lemma~\ref{thm:unique_P_R},
despite that neither $\bm{P}$ nor $\bm{R}$ is unique.
This argument is not discarded.
In the proof of injectivity of the exponential map,
the argument is present but simplified by symmetric positive definiteness.
The problem of calculating the Riemannian distance actually makes us
show a stronger result that minimizing the norm of $\bm{B}\bm{Q}-\bm{A}$
with fixed $\bm{A}$ and $\bm{B}$ will already lead to $\bm{Q} = \bm{Q}(\bm{A},\bm{B})$.
\end{remark}

\subsection{Parallel Transport}
In \cite{Sun_Lin_Liu_2026},
the parallel transport maps are always used along 
the curve minimizing Riemannian distance.
(In their work, the underlying manifold is always assumed to be 
complete and therefore such a parallel transport map is 
always well-defined.
Given that our manifolds $\Theta$ and $M$ are both not complete,
it is not immediately clear that a parallel transport map along 
distance minimizing curve should be well-defined.)
We will make the same definitions
but only for given $[\bm{A},s] \in M$ and $[\bm{B}, t] \in U_{[\bm{A},s]}$
because it is only in this case where we have shown that 
there exists a unique 
distance minimizing curve $\gamma$, given in
\eqref{eq:dist_minimizing_curve}.
We will also compute the parallel transport vector field 
and hence the parallel transport map.
(Although, in the current low dimensional document, 
the only thing we will use out of this computation 
is that the parallel transport map depends 
\textit{smoothly on the tuple of all its entries}.)
(I tried to trivialize this by searching for an existing result that proves the smoothness 
on all entries in textbooks, but no textbooks seem to have prove 
this exact statement and I have to do it in the harder way,
namely to compute everything out and see the smoothness directly.)
This is summarised in the following proposition.

\begin{proposition}\label{thm:parallel_transport}
Given $[\bm{A},s] \in M$ and $[\bm{B}, t] \in U_{[\bm{A},s]}$,
put $\bm{v} := \text{Exp}^{-1}_{[\bm{A},s]}([\bm{B}, t])$.
Then, as already proved in 
Proposition~\ref{thm:exponential_map},
with $\bm{V} = \bm{B}Q(\bm{A}, \bm{B}) - \bm{A}$
and $c = t - s$, $(\bm{V}, c)$ is the horizontal lift of $\bm{v}$ at $(\bm{A},s)$
and $\gamma: [0,1] \to M$ defined by, for all $u \in [0,1]$,

\begin{equation*}
  \gamma(u) = [\bm{A} + u \bm{V}, s + uc],
\end{equation*}
is the 
unique distance minimizing curve connecting $[\bm{A},s]$ to $[\bm{B}, t]$,
in accordance with 
\eqref{eq:dist_minimizing_curve}.
The lift of $\gamma$ to $(\bm{A},s)$ 
is the curve $L:[0,1] \to \Theta$ defined by, for all $u \in [0,1]$,

\begin{equation*}
  L(u) = (\bm{A} + u \bm{V}, s + uc).
\end{equation*}

Let us write $\Pi_{[\bm{B}, t]}^{[\bm{A},s]}: T_{[\bm{B}, t]}M \to T_{[\bm{A},s]}M$
for the parallel transport map along $\gamma$.
Given $\bm{h} \in T_{[\bm{B}, t]}M$,
let $(\bm{Z}, \beta)$ be the lift of $\bm{h}$ at $(\bm{B}\bm{Q}(\bm{A},\bm{B}), t)$
and $\bm{X}$ be the
parallel transport vector field along $\gamma$ for $\bm{h}$.
This means $\bm{X}(u) \in T_{\gamma(u)}M$ for $u \in [0,1]$,
$\bm{X}(1) = \bm{h}$, $X$ is smooth, and the covariant derivative along 
$\gamma$, denoted by $\frac{D}{du} \bm{X}(u)$, always equal to 
the zero vector field along the curve.
Then, for each $u \in [0,1]$,

\begin{equation}\label{eq:parallel_transport_X_formula}
  \bm{X}(u) = \mathrm{d} \pi_{L(u)}
  \bracket{(\bm{Z}(u), \beta)}
\end{equation}
and 
\begin{equation*}
  \Pi_{[\bm{B}, t]}^{[\bm{A},s]}(\bm{h}) = \bm{X}(0),
\end{equation*}
where $\bm{Z}(u)$, a $p \times q$ matrix, 
is the unique solution of the first order 
ordinary differential equation 

\begin{equation}\label{eq:parallel_transport_equation}
  \frac{\mathrm{d}}{\mathrm{d} u} \bm{Z}(u)
  = - \frac{\partial}{\partial w}\biggr |_{w=u} 
  (\bm{A}+w\bm{V})\psi^{-1}_{(\bm{A}+w\bm{V})^T(\bm{A}+w\bm{V})}
\bracket{(\bm{A}+w\bm{V})^T\bm{Z}(u) - \bm{Z}(u)^T(\bm{A}+w\bm{V})}
\end{equation}
subject to the initial condition $\bm{Z}(1) = \bm{Z}$.
Here $\psi$ is as in Lemma~\ref{thm:psi_S}.

Finally, as a part of the conclusion, the quantity 
$\Pi_{[\bm{B}, t]}^{[\bm{A},s]}(\bm{h})$ depends smoothly on 
the tuple $([\bm{A},s], [\bm{B}, t], \bm{h})$.
\end{proposition}

\begin{proof}
Let $\bm{Y}(u) = (\bm{Z}(u), \beta(u))$ be the horizontal lift of $\bm{X}(u)$,
which is a smooth horizontal vector field along $L$.
We prove that $\bm{Z}(u)$ will satisfy \eqref{eq:parallel_transport_equation}
and that $\beta(u) = \beta$ for all $u$
using the projection map computed in Proposition~\ref{thm:projection}.

For each point $(\bm{A},s) \in \Theta$, let us write
$P_{(\bm{A},s)}:T_{(\bm{A},s)}\Theta \to T_{(\bm{A},s)}\Theta$
for the orthogonal projection onto the horizontal tangent space 
$H_{(\bm{A},s)}\Theta$.
Given that $\bm{Y}(u)$ is a horizontal vector field, 
it is invariant under this projection, namely, for all $u \in [0,1]$:

\begin{equation}\label{eq:Y_horizontal_projection}
  \bm{Y}(u) = P_{L(u)}(\bm{Y}(u)).
\end{equation}
By O'Neill's formula (see for instance \cite{Lee_2018} equation 5.27), 
the Levi-Civita connection 
in $M$ evaluated at two vector fields in $M$ lifts hozirontally
to the Levi-Civita connection 
in $\Theta$ evaluated at the horizontally 
lifted version of the same two vector fields
plus a verticle term.
Consequently, the horizontal lift of 
$\frac{\mathrm{D}}{\mathrm{d} u} \bm{X}(u)$ 
equals to $\frac{\mathrm{D}}{\mathrm{d} u} \bm{Y}(u)$
plus a verticle term.
Given that $\bm{X}(u)$ is the parallel transport vector field,
$\frac{\mathrm{D}}{\mathrm{d} u} \bm{X}(u)$ is the zero vector field 
and its horizontal lift must also be the zero vector field.
Therefore, $\frac{\mathrm{D}}{\mathrm{d} u} \bm{Y}(u)$
is a verticle vector field and 

\begin{equation}\label{eq:Y_dot_vertical}
  P_{L(u)}(\frac{\mathrm{D}}{\mathrm{d} u} \bm{Y}(u)) = \bm{0}.
\end{equation}
With this in mind, 
covariant differentiating both sides of 
\eqref{eq:Y_horizontal_projection} to get
\begin{equation}\label{eq:covariant_euclidean}
  \frac{\mathrm{d}}{\mathrm{d} u} (\bm{Z}(u), \bm{\beta}(u))
  = \frac{\mathrm{D}}{\mathrm{d} u} \bm{Y}(u)
  = \frac{\mathrm{D}}{\mathrm{d} u} P_{L(u)}(\bm{Y}(u))
  = \frac{\mathrm{d}}{\mathrm{d} u} P_{L(u)}(\bm{Y}(u)),
\end{equation}
where we have used the fact that 
the covariant derivative in Euclidean space 
is just the ordinary derivative.  
We proceed by expanding out 
\eqref{eq:Y_dot_vertical} and \eqref{eq:covariant_euclidean}
using the chain rule, taking care of 
the dependencies of 
$P_{L(u)}(\bm{Y}(u))$ on $u$ for both entries.
Put 

\begin{equation*}
F(w, \bm{W}) := \bm{W} - (\bm{A}+w\bm{V})\psi^{-1}_{(\bm{A}+w\bm{V})^T(\bm{A}+w\bm{V})}
\bracket{(\bm{A}+w\bm{V})^T\bm{W} - \bm{W}^T(\bm{A}+w\bm{V})}.
\end{equation*}
Then $F$ is smooth by Lemma~\ref{thm:psi_S},
linear in the second argument,
and, by
\eqref{eq:horizontal_projection},
\begin{equation*}
P_{L(u)}(\bm{Y}(u)) = (F(u, \bm{Z}(u)), \beta(u)).
\end{equation*}
\eqref{eq:Y_dot_vertical} now becomes
\begin{equation}\label{eq:F_W_derivative_0}
  \bm{0}
  = P_{L(u)}\bracket{\frac{\mathrm{d}}{\mathrm{d} u}\bm{Y}(u)}
  = \bracket{F\bracket{u, \frac{\mathrm{d}}{\mathrm{d} u}\bm{Z}(u)},\
    \frac{\mathrm{d}}{\mathrm{d} u}\beta(u)}.
\end{equation}
Looking at the second component, 
we get that $\frac{\mathrm{d}}{\mathrm{d} u}\beta(u) = 0$
and that $\beta(u)$ is constant.
Next, the first component in \eqref{eq:covariant_euclidean}
becomes 
\begin{equation*}
  \begin{split}
  & \tab[0.4cm] \frac{\mathrm{d}}{\mathrm{d} u} \bm{Z}(u)
  = \frac{\mathrm{d}}{\mathrm{d} u} F(u, \bm{Z}(u))
  = \frac{\partial F}{\partial w}(u, \bm{Z}(u))
  + F\bracket{u, \frac{\mathrm{d}}{\mathrm{d} u}\bm{Z}(u)}
  \\ & = \frac{\partial F}{\partial w}(u, \bm{Z}(u))
  = - \frac{\partial}{\partial w}\biggr |_{w=u} 
  (\bm{A}+w\bm{V})\psi^{-1}_{(\bm{A}+w\bm{V})^T(\bm{A}+w\bm{V})}
\bracket{(\bm{A}+w\bm{V})^T\bm{Z}(u) - \bm{Z}(u)^T(\bm{A}+w\bm{V})}.
  \end{split}
\end{equation*}
In the second term of the second equality, 
we have used the fact that $F$ is linear in $\bm{W}$
so that its derivative is the map itself.
This second term equals to zero
by the first component of \eqref{eq:F_W_derivative_0}.

We have therefore proved \eqref{eq:parallel_transport_X_formula} 
and \eqref{eq:parallel_transport_equation}.
Given that \eqref{eq:parallel_transport_equation} 
is a system of first order differential equation with an 
initial condition with all dependencies 
being smooth,
a standard result in the theory of ordinary 
differential equations show that 
$\bm{Z}(u)$ exists, is unique, and depends smoothly 
on the tuple of everything in the equation and the initial condition.
That is, for any $u \in [0,1]$, 
$\bm{Z}(u)$ depends smoothly on 
the tuple $([\bm{A},s], [\bm{B},t], \bm{Z}, \beta)$.
Given that horizontal lifts are smooth,
$\Pi_{[\bm{B}, t]}^{[\bm{A},s]}(\bm{h})$ depends smoothly on 
the tuple $([\bm{A},s], [\bm{B}, t], \bm{h})$,
where different tangent spaces are identified together 
as the subspaces of $\mathbb{R}^{p \times q + 1}$.
This completes the proof.
\end{proof}

\subsection{A Limit Result}
We end this section with a result concerning the limit of a sequence.

\begin{proposition}\label{thm:quotient_limit}
Given a sequence 
$\{\bm{A}_n\}_{n=1}^{\infty}$
of real $p \times q$ matrices
with the property $\bm{A}_n\bm{A}_n^T \to \bm{A}\bm{A}^T$
as $n \to \infty$
for a $p \times q$ matrix $\bm{A}$
with $\text{rank}(\bm{A}) = q$, 
then 

\begin{equation*}
  \lim_{n \to \infty} \min_{\bm{Q} \in O(q)} \|\bm{A}_n\bm{Q}- \bm{A}\|_F = 0.
\end{equation*}

In particular, if $\{(\bm{A}_n, s_n)\}_{n=1}^{\infty}$
is a sequence in $\Theta$ 
such that, for some $(\bm{A}, s) \in \Theta$,
we have both $s_n \to s$ and 
$\Phi((\bm{A}_n, s_n)) 
\to \Phi((\bm{A}, s))$ as $n \to \infty$,
then $[\bm{A}_n, s_n] \to [\bm{A}, s]$ in $M$ as $n \to \infty$.
\end{proposition}

\begin{proof}
Assume, for the sake of contradiction,
that there exists an $\epsilon > 0$ and a subsequence $\{n_k\}_{k=1}^{\infty}$
such that, for all $k$,

\begin{equation*}
  \min_{\bm{Q} \in O(q)} \|\bm{A}_{n_k}\bm{Q}- \bm{A}\|_F \geqslant \epsilon.
\end{equation*}
Given that $\|\bm{A}_{n}\|_F^2 = \tr(\bm{A}_{n}\bm{A}_{n}^T)$
is convergent, in particular bounded, 
choose subsubsequence $\{n_{k_l}\}_{l=1}^{\infty}$
such that $\bm{A}_{n_{k_l}} \to \bm{B}$ as $l \to \infty$
for some $p \times q$ matrix $\bm{B}$.
It follows that $\bm{B}\bm{B}^T = \bm{A}\bm{A}^T$.
Therefore, $\text{rank}(\bm{A}) = \text{rank}(\bm{B}) = q$
and $\bm{B} = \bm{A}\bm{P}$ for some $\bm{P} \in O(q)$.
We have 
\begin{equation*}
  \min_{\bm{Q} \in O(q)} \|\bm{A}_{n_{k_l}}\bm{Q}- \bm{A}\|_F \leqslant 
  \|\bm{A}_{n_{k_l}}\bm{P}^T- \bm{A}\|_F = \|\bm{A}_{n_{k_l}}\bm{P}^T- \bm{B}\bm{P}^T\|_F
  = \|\bm{A}_{n_{k_l}}- \bm{B}\|_F \to 0
\end{equation*}
as $l \to \infty$, a contradiction.

Now suppose that $\{(\bm{A}_n, s_n)\}_{n=1}^{\infty}$
is a sequence in $\Theta$ such that, 
for some $(\bm{A}, s) \in \Theta$,
we have both $s_n \to s$ and 
$\Phi((\bm{A}_n, s_n)) 
\to \Phi((\bm{A}, s))$ as $n \to \infty$.
It follows that $\bm{A}_n\bm{A}_n^T \to \bm{A}\bm{A}^T$
and the previous paragraph shows that 
there exists a sequence $\{\bm{Q}_n\}_{n=1}^{\infty}$
of $q \times q$ orthogonal matrices 
such that $\bm{A}_n \bm{Q}_n \to \bm{A}$ 
as $n \to \infty$.
Thus,

\begin{equation*}
  \lim_{n \to \infty} \squared{\bm{A}_n, s_n}
  = \lim_{n \to \infty} \pi((\bm{A}_n\bm{Q}_n, s_n))
  = \pi((\bm{A}, s)) = \squared{\bm{A}, s}.
\end{equation*}
The proof is complete.
\end{proof}

\section{Properties of the Probabilistic Principal Component Analysis Model}\label{sec:PPCA_properties}
We now show that the statistical model $\mathcal{P}$ 
given by \eqref{eq:quotient_model}
is \textit{differentiable in quadratic mean}
and has \textit{positive definite Fisher information operator} 
in the sense developed in \cite{Sun_Lin_Liu_2026}, definition 
III.1 and III.2.

Let us introduce a few notations.
For each $[\bm{W},\sigma^2] \in M$,
let

\begin{equation}\label{eq:p_def}
  p(\bm{x}; [\bm{W},\sigma^2])
  = (2\pi)^{-\frac{p}{2}} \text{det}^{-\frac{1}{2}}(\bm{\Sigma}([\bm{W},\sigma^2])) \exp \bracket{- 
  \frac{1}{2}\bm{x}^T(\bm{\Sigma}([\bm{W},\sigma^2]))^{-1}\bm{x}}
\end{equation}
be the density of $N_p(\bm{0}, \bm{\Sigma}([\bm{W},\sigma^2]))$.
For $\bm{v} \in \mathcal{E}_{[\bm{W},\sigma^2]} \subset T_{[\bm{W},\sigma^2]}M \cong \mathbb{R}^{d}$,
write $l_{[\bm{W},\sigma^2]}(\bm{v})$
for the log-likelihood function of the local experiment
\begin{equation}\label{eq:l_def}
  l_{[\bm{W},\sigma^2]}(\bm{v})(\bm{x}) = \log 
  \bracket{p(\bm{x}; \text{Exp}_{[\bm{W},\sigma^2]}(\bm{v}))}.
\end{equation}
Define the square root likelihood function $r_{[\bm{W},\sigma^2]} : 
\mathcal{E}_{[\bm{W},\sigma^2]} \to L^2(\mathbb{R}^p)$
by setting, for all $\bm{v} \in \mathcal{E}_{[\bm{W},\sigma^2]}$,
that 
\begin{equation}\label{eq:r_def}
  r_{[\bm{W},\sigma^2]}(\bm{v})(\bm{x}) = 
  \sqrt{p(\bm{x}; \text{Exp}_{[\bm{W},\sigma^2]}(\bm{v}))}.
\end{equation}

\begin{theorem}\label{thm:DQM}
$r_{[\bm{W},\sigma^2]}$ is Fréchet differentiable at any $\bm{v} \in \mathcal{E}_{[\bm{W},\sigma^2]}$.
In particular, $\mathcal{P}$ is differentiable in quadratic mean with score
\begin{equation}\label{eq:score}
  s(\bm{x}; [\bm{W},\sigma^2])(\bm{h})
  = \frac{1}{2}\tr \curly{\bracket{\bm{\Sigma}([\bm{W},\sigma^2])}^{-1}
  \squared{\bm{x}\bm{x}^T - \bracket{\bm{\Sigma}([\bm{W},\sigma^2])}}
  \bracket{\bm{\Sigma}([\bm{W},\sigma^2])}^{-1} \bracket{D_{\bm{h}}\bm{\Sigma}([\bm{W},\sigma^2])}},
\end{equation}
where
\begin{equation}\label{eq:Sigma_directional}
  D_{\bm{h}}\bm{\Sigma}([\bm{W},\sigma^2])
  := \frac{d}{dt}\bigg|_{t=0} \bm{\Sigma}(\text{Exp}_{[\bm{W},\sigma^2]}(t\bm{h}))
  = \bm{U}\bm{W}^T + \bm{W}\bm{U}^T + b  \bm{I}_p,
\end{equation}
if $\bm{h}\in T_{[\bm{W},\sigma^2]}M$ lifts to $(\bm{U},b)$ at $(\bm{W},\sigma^2)$.
Furthermore, $\mathcal{P}$ is a regular model in the sense that
the Fisher information operator
\begin{equation}\label{eq:G_def}
  G_{[\bm{W},\sigma^2]} (\bm{h}_1,\bm{h}_2)
  = \frac{1}{2}\tr \curly{\bracket{\bm{\Sigma}([\bm{W},\sigma^2])}^{-1}
  D_{\bm{h}_1}\bm{\Sigma}([\bm{W},\sigma^2])
  \bracket{\bm{\Sigma}([\bm{W},\sigma^2])}^{-1}
  D_{\bm{h}_2}\bm{\Sigma}([\bm{W},\sigma^2])}
\end{equation}
 is positive definite.
\end{theorem}

\begin{proof}
The proof for the Fréchet differentiability
is long but routine.
We compute the Gâteaux derivative
of $r_{[\bm{W},\sigma^2]}$ at $\bm{v} \in \mathcal{E}_{[\bm{W},\sigma^2]}$
in a given direction $\bm{h} \in T_{[\bm{W},\sigma^2]}M$
and show that the Gâteaux derivative is the Fréchet derivative 
by a direct analysis of the remainder term.
At the very end, we specialize our computations to 
$\bm{v} = \bm{0}$ to prove \eqref{eq:score}
and \eqref{eq:G_def}.

Given $\bm{v}$ and $\bm{h}$,
there exists an $\epsilon > 0$ such that
$\bm{v} + t\bm{h} \in \mathcal{E}_{[\bm{W},\sigma^2]}$
whenever $t \in (-\epsilon, \epsilon)$.
If they respectively lift to $(\bm{V},c)$
and $(\bm{U},b)$ at $(\bm{W},\sigma^2)$,
then $\bm{v} + t\bm{h}$ lifts to $(\bm{V} + t\bm{U},c+tb)$
at $(\bm{W},\sigma^2)$.
First of all:

\begin{equation}\label{eq:Sigma_derivative}
  \begin{split}
  D_{\bm{h}}\bm{\Sigma}(\text{Exp}_{[\bm{W},\sigma^2]}(\bm{v})) & :=
  \frac{d}{dt}\bigg|_{t=0} \bm{\Sigma}(\text{Exp}_{[\bm{W},\sigma^2]}(\bm{v} + t\bm{h}))
  \\ & = \frac{d}{dt}\bigg|_{t=0} \curly{(\bm{W}+\bm{V} + t\bm{U})(\bm{W}+\bm{V} + t\bm{U})^T + (\sigma^2+c+tb) \bm{I}_p}
  \\ & 
  = \bm{U}\bm{V}^T + \bm{V}\bm{U}^T + \bm{U}\bm{W}^T + \bm{W}\bm{U}^T + b  \bm{I}_p,
  \end{split}
\end{equation}
proving \eqref{eq:Sigma_directional} with special case $\bm{v}=\bm{0}$.
Secondly, for all $\bm{x} \in \mathbb{R}^p$:
\begin{equation}\label{eq:l_derivative}
  \begin{split}
  & \tab[0.4cm] D_{\bm{h}}l_{[\bm{W},\sigma^2]}(\bm{v})(\bm{x}) :=
  \frac{d}{dt}\bigg|_{t=0} l_{[\bm{W},\sigma^2]}(\bm{v} + t\bm{h})(\bm{x})
  \\ & = \frac{d}{dt}\bigg|_{t=0} \curly{-\frac{p}{2}\log(2\pi) 
  -\frac{1}{2}\log \det\bracket{\bm{\Sigma}(\text{Exp}_{[\bm{W},\sigma^2]}(\bm{v} + t\bm{h}))}
  - \frac{1}{2}\bm{x}^T\bracket{\bm{\Sigma}(\text{Exp}_{[\bm{W},\sigma^2]}(\bm{v} + t\bm{h}))}^{-1}\bm{x} }
  \\ & =
  -\frac{1}{2}\tr 
  \squared{\bracket{\bm{\Sigma}(\text{Exp}_{[\bm{W},\sigma^2]}(\bm{v}))}^{-1}
  \bracket{\frac{d}{dt}\bigg|_{t=0}\bm{\Sigma}(\text{Exp}_{[\bm{W},\sigma^2]}(\bm{v} + t\bm{h})) }}
  \\ & + \frac{1}{2}\bm{x}^T\bracket{\bm{\Sigma}(\text{Exp}_{[\bm{W},\sigma^2]}(\bm{v}))}^{-1}
  \bracket{\frac{d}{dt}\bigg|_{t=0}\bm{\Sigma}(\text{Exp}_{[\bm{W},\sigma^2]}(\bm{v} + t\bm{h}))}
  \bracket{\bm{\Sigma}(\text{Exp}_{[\bm{W},\sigma^2]}(\bm{v}))}^{-1}\bm{x}
  \\ & =
  -\frac{1}{2} \tr 
  \squared{\bracket{\bm{\Sigma}(\text{Exp}_{[\bm{W},\sigma^2]}(\bm{v}))}^{-1}
  \bracket{D_{\bm{h}}\bm{\Sigma}(\text{Exp}_{[\bm{W},\sigma^2]}(\bm{v})) }}
  \\ & + \frac{1}{2}\bm{x}^T\bracket{\bm{\Sigma}(\text{Exp}_{[\bm{W},\sigma^2]}(\bm{v}))}^{-1}
  \bracket{D_{\bm{h}}\bm{\Sigma}(\text{Exp}_{[\bm{W},\sigma^2]}(\bm{v}))}
  \bracket{\bm{\Sigma}(\text{Exp}_{[\bm{W},\sigma^2]}(\bm{v}))}^{-1}\bm{x}
  \\ & = \frac{1}{2} \tr \curly{\bracket{\bm{\Sigma}(\text{Exp}_{[\bm{W},\sigma^2]}(\bm{v}))}^{-1}
  \squared{\bm{x}\bm{x}^T - \bracket{\bm{\Sigma}(\text{Exp}_{[\bm{W},\sigma^2]}(\bm{v}))}}
  \bracket{\bm{\Sigma}(\text{Exp}_{[\bm{W},\sigma^2]}(\bm{v}))}^{-1} 
  \bracket{D_{\bm{h}}\bm{\Sigma}(\text{Exp}_{[\bm{W},\sigma^2]}(\bm{v}))}}.
  \end{split}
\end{equation}
Finally:
\begin{equation}\label{eq:r_derivative}
  \begin{split}
  D_{\bm{h}}r_{[\bm{W},\sigma^2]}(\bm{v})(\bm{x})
  & = D_{\bm{h}} \exp \bracket{\frac{1}{2}l_{[\bm{W},\sigma^2]}(\bm{v})(\bm{x})}
  = \frac{1}{2} \exp \bracket{\frac{1}{2}l_{[\bm{W},\sigma^2]}(\bm{v})(\bm{x})} D_{\bm{h}}l_{[\bm{W},\sigma^2]}(\bm{v})(\bm{x})
  \\ & = \frac{1}{2} r_{[\bm{W},\sigma^2]}(\bm{v})(\bm{x}) D_{\bm{h}}l_{[\bm{W},\sigma^2]}(\bm{v})(\bm{x}).
  \end{split}
\end{equation}
To prove that \eqref{eq:r_derivative}
indeed defines a valid Gâteaux derivative,
we must prove that $D_{\bm{h}}r_{[\bm{W},\sigma^2]}(\bm{v}) \in L^2(\mathbb{R}^p)$.
Luckily, the expectation and covariance of normal 
quadratic forms have already been very well studied in 
\cite{Gupta_Nagar_2000}.
For $\bm{X} \sim N_p(\bm{0}, \bm{\Sigma})$,
Theorem 7.7.1(i) in 
\cite{Gupta_Nagar_2000}
states that 
\begin{equation*}
  \mathbb{E} \squared{\bm{X}^T \bm{Z} \bm{X}}
  = \tr \bracket{\bm{Z}\bm{\Sigma}}
\end{equation*}
for any $p \times p$ matrix $\bm{Z}$
and 
Theorem 7.7.2 in 
\cite{Gupta_Nagar_2000} states that 
\begin{equation*}
  \text{Cov} \bracket{\bm{X}^T \bm{Z} \bm{X}, \bm{X}^T \bm{Y} \bm{X}}
  = \tr(\bm{Z}\bm{\Sigma}\bm{Y}^T\bm{\Sigma}) + \tr(\bm{Z}^T\bm{\Sigma}\bm{Y}^T\bm{\Sigma})
\end{equation*} 
for any two $p \times p$ matrices $\bm{Y}$ and $\bm{Z}$.
Therefore, 
we can confirm without effort that the expectation of derivative of log density is zero:
\begin{equation*}
  \begin{split}
  & \tab[0.4cm] \int_{\mathbb{R}^p} D_{\bm{h}}l_{[\bm{W},\sigma^2]}(\bm{v})(\bm{x})
  p(\bm{x}; \bm{\Sigma}(\text{Exp}_{[\bm{W},\sigma^2]}(\bm{v})))\mathrm{d} \bm{x}
   \\ & = -\frac{1}{2} \tr 
  \squared{\bracket{\bm{\Sigma}(\text{Exp}_{[\bm{W},\sigma^2]}(\bm{v}))}^{-1}
  \bracket{D_{\bm{h}}\bm{\Sigma}(\text{Exp}_{[\bm{W},\sigma^2]}(\bm{v})) }}
  \\ & + \frac{1}{2}\int_{\mathbb{R}^p}\squared{\bm{x}^T\bracket{\bm{\Sigma}(\text{Exp}_{[\bm{W},\sigma^2]}(\bm{v}))}^{-1}
  \bracket{D_{\bm{h}}\bm{\Sigma}(\text{Exp}_{[\bm{W},\sigma^2]}(\bm{v}))}
  \bracket{\bm{\Sigma}(\text{Exp}_{[\bm{W},\sigma^2]}(\bm{v}))}^{-1}\bm{x}} 
  p(\bm{x}; \bm{\Sigma}(\text{Exp}_{[\bm{W},\sigma^2]}(\bm{v})))\mathrm{d} \bm{x} = 0
  \end{split}
\end{equation*} 
and, whenever $\bm{h}_1$ and $\bm{h}_2 \in T_{[\bm{W},\sigma^2]}M$:
\begin{equation}\label{eq:fisher_computation}
  \begin{split}
  & \tab[0.4cm] 	
   \int_{\mathbb{R}^p} D_{\bm{h}_1}l_{[\bm{W},\sigma^2]}(\bm{v})(\bm{x})
   D_{\bm{h}_2}l_{[\bm{W},\sigma^2]}(\bm{v})(\bm{x})
  p(\bm{x}; \bm{\Sigma}(\text{Exp}_{[\bm{W},\sigma^2]}(\bm{v})))\mathrm{d} \bm{x}
  \\ & = \frac{1}{4}\text{Cov}\bracket{\bm{X}^T\bm{Z}_1\bm{X}, \, \bm{X}^T\bm{Z}_2\bm{X}}
  = \frac{1}{4} \cdot 2 \tr \bracket{\bm{Z}_1 \bm{\Sigma}(\text{Exp}_{[\bm{W},\sigma^2]}(\bm{v}))
  \bm{Z}_2 \bm{\Sigma}(\text{Exp}_{[\bm{W},\sigma^2]}(\bm{v}))}
  \\ & = \frac{1}{2}\tr \curly{\bracket{\bm{\Sigma}(\text{Exp}_{[\bm{W},\sigma^2]}(\bm{v}))}^{-1}
  D_{\bm{h}_1}\bm{\Sigma}(\text{Exp}_{[\bm{W},\sigma^2]}(\bm{v}))
  \bracket{\bm{\Sigma}(\text{Exp}_{[\bm{W},\sigma^2]}(\bm{v}))}^{-1}
  D_{\bm{h}_2}\bm{\Sigma}(\text{Exp}_{[\bm{W},\sigma^2]}(\bm{v}))},
  \end{split}
\end{equation}
where $\bm{X} \sim N_p(\bm{0}, \bm{\Sigma}(\text{Exp}_{[\bm{W},\sigma^2]}(\bm{v})))$ and
\begin{equation}\label{eq:Z_j_def}
  \bm{Z}_j := \bracket{\bm{\Sigma}(\text{Exp}_{[\bm{W},\sigma^2]}(\bm{v}))}^{-1}
  \bracket{D_{\bm{h}_j}\bm{\Sigma}(\text{Exp}_{[\bm{W},\sigma^2]}(\bm{v}))}
  \bracket{\bm{\Sigma}(\text{Exp}_{[\bm{W},\sigma^2]}(\bm{v}))}^{-1}
\end{equation}
is symmetric for $j=1,2$, so that the two terms in Theorem 7.7.2
of \cite{Gupta_Nagar_2000} coincide.
Thus, in the particular case $\bm{h} = \bm{h}_1 = \bm{h}_2$:
\begin{equation}\label{eq:r_L2norm}
  \begin{split}
  & \tab[0.4cm] 	\|D_{\bm{h}}r_{[\bm{W},\sigma^2]}(\bm{v});L^2(\mathbb{R}^p)	\|^{2} =
  \int_{\mathbb{R}^p} \squared{D_{\bm{h}}r_{[\bm{W},\sigma^2]}(\bm{v})(\bm{x})}^2 \mathrm{d} \bm{x}
   \\ & = \frac{1}{4}\int_{\mathbb{R}^p} \bracket{D_{\bm{h}}l_{[\bm{W},\sigma^2]}(\bm{v})(\bm{x})}^2
  p(\bm{x}; \bm{\Sigma}(\text{Exp}_{[\bm{W},\sigma^2]}(\bm{v})))\mathrm{d} \bm{x}
  = \frac{1}{8} \tr \curly{\squared{\bracket{\bm{\Sigma}(\text{Exp}_{[\bm{W},\sigma^2]}(\bm{v}))}^{-1}
  D_{\bm{h}}\bm{\Sigma}(\text{Exp}_{[\bm{W},\sigma^2]}(\bm{v}))}^2}.
  \end{split}
\end{equation} 
Therefore, $D_{\bm{h}}r_{[\bm{W},\sigma^2]}(\bm{v}) \in L^2(\mathbb{R}^p)$
and \eqref{eq:r_derivative} is indeed a 
valid Gâteaux derivative at $\bm{v}$ in the direction $\bm{h}$.

To prove that this is indeed the Fréchet derivative,
let us look at the remainder term:

\begin{equation*}
  R(\bm{v}, \bm{h})(\bm{x}) := 
  r_{[\bm{W},\sigma^2]}(\bm{v}+\bm{h})(\bm{x}) - 
  r_{[\bm{W},\sigma^2]}(\bm{v})(\bm{x}) - D_{\bm{h}}r_{[\bm{W},\sigma^2]}(\bm{v})(\bm{x}) 
\end{equation*}
and prove that 
\begin{equation*}
  \lim_{\bm{h} \to \bm{0}} \frac{1}{\abs{\bm{h}}^2}
  \|R(\bm{v}, \bm{h});L^2(\mathbb{R}^p)	\|^2 = 0.
\end{equation*}
To start, the fundamental theorem of calculus shows that

\begin{equation*}
  \begin{split}
  & \tab[0.4cm] r_{[\bm{W},\sigma^2]}(\bm{v}+\bm{h})(\bm{x}) - 
  r_{[\bm{W},\sigma^2]}(\bm{v})(\bm{x}) 
  = \int_{0}^1 \frac{d}{dt} r_{[\bm{W},\sigma^2]}(\bm{v}+t\bm{h})(\bm{x}) \mathrm{d} t 
  \\ & = \int_{0}^1 \lim_{s \to 0} \frac{r_{[\bm{W},\sigma^2]}(\bm{v}+(t+s)\bm{h})(\bm{x})
  - r_{[\bm{W},\sigma^2]}(\bm{v}+t\bm{h})(\bm{x})}{s} \mathrm{d} t 
  = \int_{0}^1 D_{\bm{h}}r_{[\bm{W},\sigma^2]}(\bm{v}+t\bm{h})(\bm{x}) \mathrm{d} t.
  \end{split}
\end{equation*}
Therefore,
by Minkowski's integral inequality and the linearity of $D_{\bm{h}}r_{[\bm{W},\sigma^2]}(\bm{v})$
in $\bm{h}$, we have
\begin{equation*}
  \begin{split}
  & \tab[0.4cm] 
  \|R(\bm{v}, \bm{h});L^2(\mathbb{R}^p)	\|
  = \|\int_{0}^1 \bracket{D_{\bm{h}}r_{[\bm{W},\sigma^2]}(\bm{v}+t\bm{h})(\bm{x})- D_{\bm{h}}r_{[\bm{W},\sigma^2]}(\bm{v})(\bm{x})} 
  \mathrm{d} t;L^2(\mathbb{R}^p)	\|
  \\ & \leqslant \int_{0}^1 
  \|D_{\bm{h}}r_{[\bm{W},\sigma^2]}(\bm{v}+t\bm{h})- D_{\bm{h}}r_{[\bm{W},\sigma^2]}(\bm{v})
  ;L^2(\mathbb{R}^p)	\| \mathrm{d} t
  = \abs{\bm{h}} \int_{0}^1 
  \|D_{\frac{\bm{h}}{\abs{\bm{h}}}}r_{[\bm{W},\sigma^2]}(\bm{v}+t\bm{h})- D_{\frac{\bm{h}}{\abs{\bm{h}}}}r_{[\bm{W},\sigma^2]}(\bm{v})
  ;L^2(\mathbb{R}^p)	\| \mathrm{d} t
  \end{split}
\end{equation*}
and 
\begin{equation*}
  \begin{split}
  & \tab[0.4cm] \frac{1}{\abs{\bm{h}}^2}
  \|R(\bm{v}, \bm{h});L^2(\mathbb{R}^p)	\|^2 \leqslant
  \curly{\int_{0}^1 
  \|D_{\frac{\bm{h}}{\abs{\bm{h}}}}r_{[\bm{W},\sigma^2]}(\bm{v}+t\bm{h})- D_{\frac{\bm{h}}{\abs{\bm{h}}}}r_{[\bm{W},\sigma^2]}(\bm{v})
  ;L^2(\mathbb{R}^p)	\| \mathrm{d} t}^2
  \\ & \leqslant 
  \curly{\int_{0}^1 \sup_{\abs{\bm{u}} = 1}
  \|D_{\bm{u}}r_{[\bm{W},\sigma^2]}(\bm{v}+t\bm{h})- D_{\bm{u}}r_{[\bm{W},\sigma^2]}(\bm{v})
  ;L^2(\mathbb{R}^p)	\| \mathrm{d} t}^2.
  \end{split}
\end{equation*}
The proof will be complete once we can show that

\begin{equation}\label{eq:int_sup_u}
  \lim_{\bm{h} \to \bm{0}} 
  \int_{0}^1 \sup_{\abs{\bm{u}} = 1}
  \|D_{\bm{u}}r_{[\bm{W},\sigma^2]}(\bm{v}+t\bm{h})- D_{\bm{u}}r_{[\bm{W},\sigma^2]}(\bm{v})
  ;L^2(\mathbb{R}^p)	\| \mathrm{d} t = 0.
\end{equation}
We must therefore examine the continuity of $D_{\bm{h}}r_{[\bm{W},\sigma^2]}(\bm{v})$,
which is easy to see if we revisit the formulas.
\eqref{eq:Sigma_derivative} shows that $ D_{\bm{h}}\bm{\Sigma}(\text{Exp}_{[\bm{W},\sigma^2]}(\bm{v}))$
is smooth in both $\bm{h}$ and $\bm{v}$.
Thus,
\eqref{eq:r_derivative} shows that,
for each fixed $\bm{x} \in \mathbb{R}^p$,
$D_{\bm{h}}r_{[\bm{W},\sigma^2]}(\bm{v})(\bm{x})$ 
is smooth in both $\bm{h}$ and $\bm{v}$
and \eqref{eq:r_L2norm}
shows that $\|D_{\bm{h}}r_{[\bm{W},\sigma^2]}(\bm{v});L^2(\mathbb{R}^p)	\|$
is smooth in both $\bm{h}$ and $\bm{v}$.
Scheffé's lemma shows that
$D_{\bm{h}}r_{[\bm{W},\sigma^2]}(\bm{v})$, as a map to $L^2(\mathbb{R}^p)$,
is continuous in $\bm{h}$ and in $\bm{v}$.
The integrand inside \eqref{eq:int_sup_u}
is thus bounded uniformly 
over all $t \in [0,1]$ and all $\bm{h}$
in a compact neighborhood of $\bm{0}$
and it suffices to show 
that, for every $t \in [0,1]$, that 
\begin{equation}\label{eq:lim_sup_u}
  \lim_{\bm{h} \to \bm{0}} 
  \sup_{\abs{\bm{u}} = 1}
  \|D_{\bm{u}}r_{[\bm{W},\sigma^2]}(\bm{v}+t\bm{h})- D_{\bm{u}}r_{[\bm{W},\sigma^2]}(\bm{v})
  ;L^2(\mathbb{R}^p)	\|  = 0,
\end{equation}
thanks to the bounded convergence theorem.
To prove \eqref{eq:lim_sup_u}, 
let $\bm{e}_1, \ldots, \bm{e}_d$ 
be a $g_{[\bm{W},\sigma^2]}$ orthonormal basis.
Every unit vector $\bm{u}$ can be 
expressed as $\bm{u} = \sum_{i=1}^d u_i \bm{e}_i$
with $\sum_{i=1}^d u_i^2 = 1$.
By linearity of $D_{\bm{h}}r_{[\bm{W},\sigma^2]}(\bm{v})$ in $\bm{h}$, 
we have 
\begin{equation*}
  \begin{split}
  & \tab[0.4cm] \|D_{\bm{u}}r_{[\bm{W},\sigma^2]}(\bm{v}+t\bm{h}) - D_{\bm{u}}r_{[\bm{W},\sigma^2]}(\bm{v});L^2(\mathbb{R}^p)	\| 
  \leqslant \sum_{i=1}^d \abs{u_i}\|D_{\bm{e}_i}r_{[\bm{W},\sigma^2]}(\bm{v}+t\bm{h}) - 
  D_{\bm{e}_i}r_{[\bm{W},\sigma^2]}(\bm{v});L^2(\mathbb{R}^p)	\|
  \\ & \leqslant \sqrt{d} \max_{i=1, \ldots, d} \|D_{\bm{e}_i}r_{[\bm{W},\sigma^2]}(\bm{v}+t\bm{h}) - 
  D_{\bm{e}_i}r_{[\bm{W},\sigma^2]}(\bm{v});L^2(\mathbb{R}^p)	\|.
  \end{split}
\end{equation*}
Taking supremum yields
\begin{equation*}
  \sup_{\abs{\bm{u}} = 1}
  \|D_{\bm{u}}r_{[\bm{W},\sigma^2]}(\bm{v}+t\bm{h})- D_{\bm{u}}r_{[\bm{W},\sigma^2]}(\bm{v})
  ;L^2(\mathbb{R}^p)	\|  \leqslant 
  \sqrt{d} \max_{i=1, \ldots, d} \|D_{\bm{e}_i}r_{[\bm{W},\sigma^2]}(\bm{v}+t\bm{h}) - 
  D_{\bm{e}_i}r_{[\bm{W},\sigma^2]}(\bm{v});L^2(\mathbb{R}^p)	\|.
\end{equation*}
Given that $D_{\bm{e}_i}r_{[\bm{W},\sigma^2]}(\bm{v}+t\bm{h}) 
\to D_{\bm{e}_i}r_{[\bm{W},\sigma^2]}(\bm{v})$
in $L^2(\mathbb{R}^p)$
as $\bm{h} \to \bm{0}$
for every $i$, every $t$, and every $\bm{v}$,
the proof of Fréchet differentiability is complete.

Finally,
the formulas \eqref{eq:Sigma_derivative}, \eqref{eq:l_derivative},
and \eqref{eq:r_derivative}
simplifies when $\bm{v} = \bm{0}$:
\begin{equation*}
  \begin{split}
  & \tab[0.4cm] D_{\bm{h}}r_{[\bm{W},\sigma^2]}(\bm{0})(\bm{x})
  = \frac{1}{2} r_{[\bm{W},\sigma^2]}(\bm{0})(\bm{x}) D_{\bm{h}}l_{[\bm{W},\sigma^2]}(\bm{0})(\bm{x})
  \\ & = \frac{1}{4} \sqrt{p(\bm{x}; [\bm{W},\sigma^2])}
  \tr \curly{\bracket{\bm{\Sigma}([\bm{W},\sigma^2])}^{-1}
  \squared{\bm{x}\bm{x}^T - \bracket{\bm{\Sigma}([\bm{W},\sigma^2])}}
  \bracket{\bm{\Sigma}([\bm{W},\sigma^2])}^{-1} \bracket{D_{\bm{h}}\bm{\Sigma}([\bm{W},\sigma^2])}}.
  \end{split}
\end{equation*}
The score $s(\bm{x}; [\bm{W},\sigma^2])(\bm{h})$
therefore equals to $D_{\bm{h}}l_{[\bm{W},\sigma^2]}(\bm{0})(\bm{x})$
and is indeed given by \eqref{eq:score}.
The Fisher information operator \eqref{eq:G_def} is obtained 
by taking $\bm{v} = \bm{0}$ in \eqref{eq:fisher_computation}.
To show that $G_{[\bm{W},\sigma^2]}$ is positive definite,
let $\bm{h} \in T_{[\bm{W},\sigma^2]}M$ be given.
We show that $G_{[\bm{W},\sigma^2]}(\bm{h},\bm{h}) \geqslant 0$
and equality holds only when $\bm{h} = \bm{0}$.
The non-negative part is clear from \eqref{eq:G_def}:

\begin{equation*}
  G_{[\bm{W},\sigma^2]}(\bm{h},\bm{h}) = 
  \frac{1}{2}\tr \curly{ \squared{\bracket{\bm{\Sigma}([\bm{W},\sigma^2])}^{-1}
  D_{\bm{h}}\bm{\Sigma}([\bm{W},\sigma^2])}^2 } \geqslant 0.
\end{equation*}
Assume from now on that the equality holds.
A slight annoyance here is that the product of
two symmetric matrices is not necessarily symmetric
and we cannot immediately conclude that 
$D_{\bm{h}}\bm{\Sigma}([\bm{W},\sigma^2]) = \bm{0}$.
We easily overcome this difficulty 
by using the cyclic property of the trace 
to get a symmetric matrix 

\begin{equation*}
  \bm{N} := \bracket{\bm{\Sigma}([\bm{W},\sigma^2])}^{-\frac{1}{2}}
  D_{\bm{h}}\bm{\Sigma}([\bm{W},\sigma^2]) 
  \bracket{\bm{\Sigma}([\bm{W},\sigma^2])}^{-\frac{1}{2}}
\end{equation*}
such that 

\begin{equation*}
  0=G_{[\bm{W},\sigma^2]}(\bm{h},\bm{h}) = 
  \frac{1}{2}\tr(\bm{N}^2).
\end{equation*}
Therefore, $\bm{N} = \bm{0}$
and $D_{\bm{h}}\bm{\Sigma}([\bm{W},\sigma^2]) = \bm{0}$.
Now \eqref{eq:Sigma_directional}
reads the same as \eqref{eq:ker_Phi_equation}
in the proof of Proposition~\ref{thm:horizontal_vertical}.
We therefore get $(\bm{U},b) \in \ker(d\Phi_{(\bm{W},\sigma^2)}) = V_{(\bm{W},\sigma^2)}$.
However, being lifted from $\bm{h}$,
$(\bm{U},b)$ must also be a horizontal vector.
Therefore, $\bm{h} = \bm{0}$.
We have proven that $G_{[\bm{W},\sigma^2]}$ is positive definite.
\end{proof}

As an immediate consequence, we have  
local asymptotic normality.

\begin{proposition}\label{thm:LAN}
With the true parameter $[\theta_0] \in M$,
let us write

\begin{equation}\label{eq:theta_n_h}
  \squared{\theta_{n, \bm{h}}} = \text{Exp}_{[\theta_0]} \bracket{\frac{\bm{h}}{\sqrt{n}}},
\end{equation}
whenever $n$ is large enough such that $\frac{\bm{h}}{\sqrt{n}} \in U_{[\theta_0]}$.
For any $[\theta] \in M$, 
let us write $P^n_{[\theta]}$
for the joint distribution of $(\bm{X}_1, \ldots, \bm{X}_n)$
when all components are independent and 
identically distributed as $N_p(\bm{0}, \Sigma([\theta]))$.

We have the local asymptotic normality expansion

\begin{equation}\label{eq:LAN_prob_convergence}
  \sum_{i=1}^n \log \frac{p\bracket{\bm{X}_i; \squared{\theta_{n, \bm{h}}}}}{p(\bm{X}_i; \theta_0)}
  - \sum_{i=1}^n s(\bm{X}_i;[\theta_0])\bracket{\frac{\bm{h}}{\sqrt{n}}}
  + \frac{1}{2} G_{[\theta_0]}(\bm{h},\bm{h}) \to 0,
\end{equation}
in probability as $n \to \infty$ when 
$(\bm{X}_1, \ldots, \bm{X}_n)$ follows $P^n_{\theta_0}$.

In particular,
\begin{equation}\label{eq:LAN_weak_convergence}
  \sum_{i=1}^n \log \frac{p\bracket{\bm{X}_i; \squared{\theta_{n, \bm{h}}}}}{p(\bm{X}_i; [\theta_0])}
  \to N \bracket{-\frac{1}{2}G_{[\theta_0]}(\bm{h},\bm{h}), \, G_{[\theta_0]}(\bm{h},\bm{h})}
\end{equation}
in distribution as $n \to \infty$
when $(\bm{X}_1, \ldots, \bm{X}_n)$ follows $P^n_{[\theta_0]}$
and
Le Cam's first lemma in the format of Example 6.5 in \cite{van_der_Vaart_1998}
shows that $P^n_{\squared{\theta_0}}$
and $P^n_{\squared{\theta_{n, \bm{h}}}}$ are mutually contiguous.
\end{proposition}

\begin{proof}
Given that $\mathcal{P}$ is differentiable in quadratic mean
and has positive definite Fisher operator,
this follows from Proposition III.1 in \cite{Sun_Lin_Liu_2026}.
\end{proof}

\section{Properties of the Maximum Likelihood Estimator of $\mathcal{P}$}\label{sec:MLE}

\subsection{The Maximum Likelihood Estimator}
In this section, we assume that 
there is a true parameter $[\bm{W}_0, \sigma_0^2] \in M$
for the model $\mathcal{P}$ given in \eqref{eq:quotient_model}
and we investigate the consistency and asymptotic normality
of the maximum likelihood estimator towards this true parameter.

More concretely,
given data points $\bm{x}_1, \ldots, \bm{x}_n \in \mathbb{R}^p$,
we consider them as realizations of independent 
identically distributed random 
variables $\bm{X}_1, \ldots, \bm{X}_n$
whose true density is $p(\bm{x}; [\bm{W}_0, \sigma_0^2])$,
where $p$ is given in \eqref{eq:p_def}.
In what follows, 
we denote the true covariance 
matrix by 
\begin{equation}\label{eq:Sigma_0_def}
  \bm{\Sigma}_0 := \bm{\Sigma}([\bm{W}_0, \sigma_0^2])
= \bm{W}_0\bm{W}_0^T + \sigma_0^2\bm{I}_p.
\end{equation}
The eigenvalues of the true covariance matrix are

\begin{equation}\label{eq:lambda_j_Sigma_0}
  \lambda_j(\bm{\Sigma}_0)
  = 
  \begin{cases}
    \sigma_j(\bm{W}_0)^2 + \sigma_0^2 \tab[4.0cm] j = 1, \ldots, q \\
    \sigma_0^2                        \tab[6.0cm] j = q+1, \ldots, p.
  \end{cases}
\end{equation}
The sample covariance matrix is defined by 

\begin{equation}\label{eq:S_n_def}
  \bm{S}_n = \frac{1}{n} \sum_{i=1}^n \bm{x}_i\bm{x}_i^T.
\end{equation}
Given that we will discuss $\bm{S}_n$ a lot, 
we shall write $\lambda_j$ in place of $\lambda_j(\bm{S}_n)$.
We need also to fix one choice of 
orthonormal basis of eigenvectors $\bm{v}_1, \ldots, \bm{v}_p$
where $\bm{v}_j = \bm{v}_j(\bm{S}_n)$ corresponds to the eigenvalue $\lambda_j$.
The facts we will use 
repeatedly is that, when $p$ and $q$ are both fixed,
$\bm{S}_n \to \bm{\Sigma}_0$ almost surely as $n \to \infty$
and, due to the fact that eigenvalues depends continuously 
on the matrices, $\lambda_j \to \lambda_j(\bm{\Sigma}_0)$
almost surely as $n \to \infty$.

We now formulate the maximum likelihood estimator as 
computed 
in \cite{Tipping_Bishop_1999}.
Put $\bm{U} = [\bm{v}_1, \ldots, \bm{v}_q]$, a $p \times q$ matrix, and 
$\Delta_q$ to be a $q \times q$ diagonal matrix with diagonal entries $\lambda_1,
\ldots, \lambda_q$.
Then, the collection of all choices 
of $(\hat{\bm{W}}_n, \hat{\sigma}_n^2)$ maximizing the likelihood 
is given by

\begin{equation}\label{eq:W_hat_def}
  \hat{\bm{W}}_n = \bm{U} \bracket{\Delta_q - \hat{\sigma}_n^2\bm{I}_q}^{\frac{1}{2}}\bm{R}
\end{equation}
and 
\begin{equation}\label{eq:sigma_hat_def}
  \hat{\sigma}_n^2 = \frac{1}{p-q} \sum_{j=q+1}^p\lambda_j,
\end{equation}
ranged over all rotation matrices $\bm{R} \in O(q)$.

It is immediately clear from \eqref{eq:lambda_j_Sigma_0}
that, almost surely, for all sufficiently large $n$,
$\lambda_j > 0$ for all $j$ and $\hat{\sigma}_n^2 > 0$
so that $\hat{\sigma}_n^2$ is well-defined.
Also, $\hat{\sigma}_n^2 \to \sigma_0^2$
almost surely as $n \to \infty$.
Another look at \eqref{eq:lambda_j_Sigma_0}
shows that, almost surely, for all sufficiently large $n$, 
$\lambda_j - \hat{\sigma}_n^2 > 0$
for all $j \in \{1, \ldots, q\}$,
$\lambda_q > \lambda_{q+1}$,
and $\text{rank}(\hat{\bm{W}}_n) = q$
so that $\hat{\bm{W}}_n$ is well-defined.
Thus, almost surely, for all sufficiently large $n$, 
the same model parametrized by $M$
has a unique maximum likelihood estimator 
$\squared{\hat{\bm{W}}_n,\hat{\sigma}_n^2}$.

\subsection{Consistency}
First of all, I would like to argue that 
this quotient manifold parameter space allows 
us to settle consistency in an (arguably more straightforward) 
way that does not use
\cite{Datta_Chakrabarty_2023}
or \cite{Redner_1981}.

\begin{theorem}\label{thm:consistency}
The maximum likelihood estimator is consistent.
That is,
$\squared{\hat{\bm{W}}_n,\hat{\sigma}_n^2} \to \squared{\bm{W}_0,\sigma_0^2}$
almost surely as $n \to \infty$.
\end{theorem}
\begin{proof}
Put

\begin{equation}\label{eq:Sigma_hat_def}
  \hat{\bm{\Sigma}}_n := \bm{\Sigma}\bracket{ \squared{\hat{\bm{W}}_n,\hat{\sigma}_n^2}}
  = \hat{\bm{W}}_n\hat{\bm{W}}_n^T + \hat{\sigma}_n^2\bm{I}_p.
\end{equation}
In light of 
Proposition~\ref{thm:quotient_limit}
and the fact that $\hat{\sigma}_n^2$ is consistent,
it suffices to prove that 
$\hat{\bm{\Sigma}}_n \to \bm{\Sigma}_0$
almost surely as $n \to \infty$.
The fact that $\sum_{j=1}^p \bm{v}_j\bm{v}_j^T = \bm{I}_p$
allows us to write

\begin{equation*}
\hat{\bm{\Sigma}}_n =   \bm{U} \bracket{\Delta_q - \hat{\sigma}_n^2\bm{I}_q}\bm{U}^T + \hat{\sigma}_n^2 \bm{I}_p
= \sum_{j=1}^q (\lambda_j-\hat{\sigma}_n^2)\bm{v}_j\bm{v}_j^T + \hat{\sigma}_n^2 \bm{I}_p
= \sum_{j=1}^q \lambda_j\bm{v}_j\bm{v}_j^T + \hat{\sigma}_n^2 \sum_{j=q+1}^p \bm{v}_j\bm{v}_j^T
\end{equation*}
and 
\begin{equation*}
\bm{S}_n = \bm{S}_n \sum_{j=1}^p \bm{v}_j\bm{v}_j^T = \sum_{j=1}^p \lambda_j\bm{v}_j\bm{v}_j^T.
\end{equation*}
Thus, 
\begin{equation*}
\hat{\bm{\Sigma}}_n-\bm{S}_n = \sum_{j=q+1}^p \bracket{\hat{\sigma}_n^2-\lambda_j} \bm{v}_j\bm{v}_j^T
\end{equation*}
and 
\begin{equation*}
	\|\hat{\bm{\Sigma}}_n-\bm{S}_n	\| \leqslant \sum_{j=q+1}^p \abs{\hat{\sigma}_n^2-\lambda_j} \|\bm{v}_j\bm{v}_j^T\| = 
  \sum_{j=q+1}^p \abs{\hat{\sigma}_n^2-\lambda_j}.
\end{equation*}
Therefore, given that both $\hat{\sigma}_n^2$ and $\lambda_j$ 
will converge almost surely to $\sigma_0^2$ when $j \in \{q+1,\ldots,p\}$,
 $\|\hat{\bm{\Sigma}}_n-\bm{S}_n	\| \to 0 $ almost surely as $n \to \infty$.
The proof is complete.
\end{proof}

\subsection{Asymptotic Linearity}

We now move to prove that,
locally at the true parameter, 
the maximum likelihood estimator 
is asymptotically normal in the tangent space 
of the true parameter by local linearity.
Given that the manifold is locally Euclidean,
this result can be attacked 
with 
the classical approach,
namely by verifying conditions of 
Theorem 5.23 in \cite{van_der_Vaart_1998},
without too much 
use of geometric analysis.
We start with computational lemmas.

\begin{lemma}
Given $[\bm{W}, \sigma^2] \in M$
and any
convex neighborhood $U$ of $\bm{0}$
in $T_{[\bm{W}, \sigma^2]}M$ 
such that $\overline{U} \subset \mathcal{E}_{[\bm{W}, \sigma^2]}$,
there exists a function $m: \mathbb{R}^p \to \mathbb{R}$
such that, whenever 
$\bm{v}_1$, $\bm{v}_2 \in U$
and $\bm{x} \in \mathbb{R}^p$, we always have 

\begin{equation*}
  \abs{l_{[\bm{W}, \sigma^2]}(\bm{v}_1)(\bm{x})-
  l_{[\bm{W}, \sigma^2]}(\bm{v}_2)(\bm{x})} 
  \leqslant m(\bm{x}) \abs{\bm{v}_1-\bm{v}_2}.
\end{equation*}
This $m$ has its second moment 
under any distribution $\mu$ on $\mathbb{R}^p$ that 
has the fourth moment.
That is,

\begin{equation*}
  \int_{\mathbb{R}^p} (m(\bm{x}))^2 \mathrm{d} \mu(\bm{x}) < \infty
\end{equation*}
whenever 
\begin{equation*}
  \int_{\mathbb{R}^p} \abs{\bm{x}}^4 \mathrm{d} \mu(\bm{x}) < \infty.
\end{equation*}
\end{lemma}
\begin{proof}
Given $[\bm{W}, \sigma^2] \in M$
and a convex neighborhood $U$ of $\bm{0}$ 
compactly contained in $\mathcal{E}_{[\bm{W}, \sigma^2]}$,
the result will follow by the mean value inequality 
if we can find a function $m$ such that 

\begin{equation*}
  \abs{D_{\bm{h}}l_{[\bm{W}, \sigma^2]}(\bm{v})(\bm{x})} \leqslant m(\bm{x})
\end{equation*}
for all unit vector $\bm{h} \in T_{[\bm{W}, \sigma^2]}M$,
all $\bm{v} \in \overline{U}$,
and all $\bm{x} \in \mathbb{R}^p$.
Given that we have proved in Proposition~\ref{thm:DQM}
that $\bracket{\bm{\Sigma}(\text{Exp}_{[\bm{W},\sigma^2]}(\bm{v}))}^{-1}$
is continuous in $\bm{v}$ 
and that $D_{\bm{h}}\bm{\Sigma}(\text{Exp}_{[\bm{W},\sigma^2]}(\bm{v}))$
is continuous in both $\bm{h}$ and $\bm{v}$,
we choose constants $C_1 \in (0, \infty)$ and $C_2 \in (0, \infty)$
such that 

\begin{equation*}
  \frac{1}{2}\left \|\bracket{\bm{\Sigma}(\text{Exp}_{[\bm{W},\sigma^2]}(\bm{v}))}^{-1}
  \bracket{D_{\bm{h}}\bm{\Sigma}(\text{Exp}_{[\bm{W},\sigma^2]}(\bm{v}))}
  \bracket{\bm{\Sigma}(\text{Exp}_{[\bm{W},\sigma^2]}(\bm{v}))}^{-1} \right \|_F
  \leqslant C_1
\end{equation*}
and 
\begin{equation*}
  \frac{1}{2}\left \|\bracket{\bm{\Sigma}(\text{Exp}_{[\bm{W},\sigma^2]}(\bm{v}))}^{-1}
  \bracket{D_{\bm{h}}\bm{\Sigma}(\text{Exp}_{[\bm{W},\sigma^2]}(\bm{v}))} \right \|_F
  \leqslant C_2
\end{equation*}
for all all unit vector $\bm{h} \in T_{[\bm{W}, \sigma^2]}M$ and
all $\bm{v} \in \overline{U}$,
thanks to the fact that $\overline{U}$ is compact.
from Proposition~\ref{thm:exponential_map}.
By \eqref{eq:l_derivative},
we have 

\begin{equation*}
  \begin{split}
  \abs{D_{\bm{h}}l_{[\bm{W}, \sigma^2]}(\bm{v})(\bm{x})} & 
  \leqslant \frac{1}{2} 	\|\bm{x}\bm{x}^T\|_F 
  \left \|\bracket{\bm{\Sigma}(\text{Exp}_{[\bm{W},\sigma^2]}(\bm{v}))}^{-1}
  \bracket{D_{\bm{h}}\bm{\Sigma}(\text{Exp}_{[\bm{W},\sigma^2]}(\bm{v}))}
  \bracket{\bm{\Sigma}(\text{Exp}_{[\bm{W},\sigma^2]}(\bm{v}))}^{-1} \right \|_F
  \\ & + \frac{1}{2}
  \left \|\bracket{\bm{\Sigma}(\text{Exp}_{[\bm{W},\sigma^2]}(\bm{v}))}^{-1}
  \bracket{D_{\bm{h}}\bm{\Sigma}(\text{Exp}_{[\bm{W},\sigma^2]}(\bm{v}))} \right \|_F
  \leqslant C_1 	\bm{x}^T\bm{x}
  + C_2
  \end{split}
\end{equation*}
Setting $m(\bm{x}) = C_1 	\bm{x}^T\bm{x}+ C_2$ will finish the proof.
\end{proof}

\begin{lemma}
Define function $L:\mathcal{E}_{[\bm{W}_0, \sigma_0^2]}\to \mathbb{R}$
by setting, for all $\bm{v} \in \mathcal{E}_{[\bm{W}_0, \sigma_0^2]}$,

\begin{equation*}
  L(\bm{v}) = \int_{\mathbb{R}^p} l_{[\bm{W}_0, \sigma_0^2]}(\bm{v})(\bm{x})
  p(\bm{x}; \bm{\Sigma}_0) \mathrm{d} \bm{x}.
\end{equation*}
Then $\bm{0}$ is the unique maximizer of $L$ and,
for all $\bm{v} \in \mathcal{E}_{[\bm{W}_0, \sigma_0^2]}$:

\begin{equation*}
  L(\bm{v}) = L(\bm{0}) - \frac{1}{2} G_{[\bm{W}_0, \sigma_0^2]}(\bm{v},\bm{v}) + R(\bm{v})
\end{equation*}
with the remainder of the second order Taylor expansion satisfying

\begin{equation*}
  \lim_{\abs{\bm{v}} \to 0} \frac{R(\bm{v})}{\abs{\bm{v}}^2} = 0.
\end{equation*}
\end{lemma}

\begin{proof}
One should immediately recognize that 
$L(\bm{0})-L(\bm{v})$ equals to the Kullback-Leibler divergence of the true model 
$N_p(\bm{0}, \bm{\Sigma}_0)$
with respect to the 
model $N_p \bracket{\bm{0}, \bm{\Sigma} \bracket{\text{Exp}_{[\bm{W}_0, \sigma_0^2]}(\bm{v})}}$.
Therefore, $L(\bm{v})$ is maximized at all $\bm{v}$
with $\bm{\Sigma} \bracket{\text{Exp}_{[\bm{W}_0, \sigma_0^2]}(\bm{v})} = \bm{\Sigma}_0$.
Given that $\bm{\Sigma}$ is injective and that we have proven 
in Proposition~\ref{thm:exponential_map} that
$\text{Exp}_{[\bm{W}_0, \sigma_0^2]}$ is injective in $\mathcal{E}_{[\bm{W}_0, \sigma_0^2]}$,
this only happens when $\bm{v} = \bm{0}$ and $\bm{0}$ is the unique maximizer of $L$.

To obtain the Taylor expansion, we basically have to repeat 
the computation part in the proof of Proposition~\ref{thm:DQM} on a slightly 
different problem with the second derivative at $\bm{0}$ being computed as well.
The first order derivative at $\bm{0}$ is automatically $0$ because $\bm{0}$ is the maximizer.

Let $\bm{h} \in T_{[\bm{W}_0, \sigma_0^2]}M$ be given.
We can easily compute that

\begin{equation*}
  L(t \bm{h}) = - \frac{p}{2} \log(2\pi) - \frac{1}{2}\log\bracket{\det \bracket{\bm{\Sigma}\bracket{\text{Exp}_{[\bm{W}_0, \sigma_0^2]}(t\bm{h})}}}
  - \frac{1}{2} \tr\curly{
  \bracket{\bm{\Sigma}\bracket{\text{Exp}_{[\bm{W}_0, \sigma_0^2]}(t\bm{h})}}^{-1} \bm{\Sigma}_0}.
\end{equation*}
Suppose that $\bm{h}$ lifts to $(\bm{U},b)$ at $(\bm{W}_0,\sigma_0^2)$.
\eqref{eq:Sigma_derivative} immediately shows that
\begin{equation*}
  \frac{d^2}{dt^2}\bigg|_{t=0} \bracket{\bm{\Sigma}\bracket{\text{Exp}_{[\bm{W}_0, \sigma_0^2]}(t\bm{h})}}
  = 2 \bm{U}\bm{U}^T.
\end{equation*}
We now compute the second derivative $L$ term by term.
For the log determinant term, we have

\begin{equation*}
  \frac{d}{dt} \bracket{- \frac{1}{2} \log\bracket{\det \bracket{\bm{\Sigma}\bracket{\text{Exp}_{[\bm{W}_0, \sigma_0^2]}(t\bm{h})}}}}
  = - \frac{1}{2} \tr \squared{\bracket{\bm{\Sigma}\bracket{\text{Exp}_{[\bm{W}_0, \sigma_0^2]}(t\bm{h})}}^{-1}
  \frac{d}{dt} \bracket{\bm{\Sigma}\bracket{\text{Exp}_{[\bm{W}_0, \sigma_0^2]}(t\bm{h})}}},
\end{equation*}
\begin{equation*}
  \begin{split}
  & \tab[0.4cm]\frac{d^2}{dt^2} \bracket{- \frac{1}{2} \log\bracket{\det \bracket{\bm{\Sigma}\bracket{\text{Exp}_{[\bm{W}_0, \sigma_0^2]}(t\bm{h})}}}}
  = - \frac{1}{2}  \frac{d}{dt} \tr \squared{\bracket{\bm{\Sigma}\bracket{\text{Exp}_{[\bm{W}_0, \sigma_0^2]}(t\bm{h})}}^{-1}
  \frac{d}{dt} \bracket{\bm{\Sigma}\bracket{\text{Exp}_{[\bm{W}_0, \sigma_0^2]}(t\bm{h})}}}
  \\ & = - \frac{1}{2} \tr \curly{\bracket{\bm{\Sigma}\bracket{\text{Exp}_{[\bm{W}_0, \sigma_0^2]}(t\bm{h})}}^{-1}
  \frac{d^2}{dt^2} \bracket{\bm{\Sigma}\bracket{\text{Exp}_{[\bm{W}_0, \sigma_0^2]}(t\bm{h})}}}
  \\ & + \frac{1}{2} \tr \curly{\squared{\bracket{\bm{\Sigma}\bracket{\text{Exp}_{[\bm{W}_0, \sigma_0^2]}(t\bm{h})}}^{-1}
  \frac{d}{dt} \bracket{\bm{\Sigma}\bracket{\text{Exp}_{[\bm{W}_0, \sigma_0^2]}(t\bm{h})}}}^2 },
  \end{split}
\end{equation*}
and therefore 
\begin{equation*}
\frac{d^2}{dt^2}\bigg|_{t=0} 
  \bracket{- \frac{1}{2} \log\bracket{\det \bracket{\bm{\Sigma}\bracket{\text{Exp}_{[\bm{W}_0, \sigma_0^2]}(t\bm{h})}}}}
  = -  \tr \curly{\bm{\Sigma}_0^{-1}\bm{U}\bm{U}^T}
  + \frac{1}{2} \tr \curly{ \squared{\bm{\Sigma}_0^{-1}
  D_{\bm{h}} \bm{\Sigma}([\bm{W}_0, \sigma_0^2])}^2}.
\end{equation*}
Next, for the trace term, we have
\begin{equation*}
  \begin{split}
  & \tab[0.4cm] \frac{d}{dt} \bracket{- \frac{1}{2} \tr\curly{
  \bracket{\bm{\Sigma}\bracket{\text{Exp}_{[\bm{W}_0, \sigma_0^2]}(t\bm{h})}}^{-1} \bm{\Sigma}_0}}
  \\ & = \frac{1}{2} \tr \curly{\bracket{\bm{\Sigma}\bracket{\text{Exp}_{[\bm{W}_0, \sigma_0^2]}(t\bm{h})}}^{-1}
  \frac{d}{dt} \bracket{\bm{\Sigma}\bracket{\text{Exp}_{[\bm{W}_0, \sigma_0^2]}(t\bm{h})}}
  \bracket{\bm{\Sigma}\bracket{\text{Exp}_{[\bm{W}_0, \sigma_0^2]}(t\bm{h})}}^{-1}\bm{\Sigma}_0},
  \end{split}
\end{equation*}
\begin{equation*}
  \begin{split}
  & \tab[0.4cm] \frac{d^2}{dt^2} \bracket{- \frac{1}{2} \tr\curly{
  \bracket{\bm{\Sigma}\bracket{\text{Exp}_{[\bm{W}_0, \sigma_0^2]}(t\bm{h})}}^{-1} \bm{\Sigma}_0}}
  \\ & = 
  \frac{1}{2} \tr 
  \biggr \{-2
  \biggr \{ \bracket{\bm{\Sigma}\bracket{\text{Exp}_{[\bm{W}_0, \sigma_0^2]}(t\bm{h})}}^{-1}
  \frac{d}{dt} \bracket{\bm{\Sigma}\bracket{\text{Exp}_{[\bm{W}_0, \sigma_0^2]}(t\bm{h})}}
  \bracket{\bm{\Sigma}\bracket{\text{Exp}_{[\bm{W}_0, \sigma_0^2]}(t\bm{h})}}^{-1}
  \\ & \frac{d}{dt} \bracket{\bm{\Sigma}\bracket{\text{Exp}_{[\bm{W}_0, \sigma_0^2]}(t\bm{h})}}
  \bracket{\bm{\Sigma}\bracket{\text{Exp}_{[\bm{W}_0, \sigma_0^2]}(t\bm{h})}}^{-1}\bm{\Sigma}_0
  \biggr \} \biggr \} \\ & +
  \frac{1}{2} \tr 
  \curly{\bracket{\bm{\Sigma}\bracket{\text{Exp}_{[\bm{W}_0, \sigma_0^2]}(t\bm{h})}}^{-1}
  \frac{d^2}{dt^2} \bracket{\bm{\Sigma}\bracket{\text{Exp}_{[\bm{W}_0, \sigma_0^2]}(t\bm{h})}}
  \bracket{\bm{\Sigma}\bracket{\text{Exp}_{[\bm{W}_0, \sigma_0^2]}(t\bm{h})}}^{-1}\bm{\Sigma}_0},
  \end{split}
\end{equation*}
and therefore
\begin{equation*}
  \frac{d^2}{dt^2}\bigg|_{t=0} \bracket{- \frac{1}{2} \tr\curly{
  \bracket{\bm{\Sigma}\bracket{\text{Exp}_{[\bm{W}_0, \sigma_0^2]}(t\bm{h})}}^{-1} \bm{\Sigma}_0}}
  = - \tr \curly{\squared{\bm{\Sigma}_0^{-1}D_{\bm{h}} \bm{\Sigma}([\bm{W}_0, \sigma_0^2])}^2}
  + \tr \bracket{\bm{\Sigma}_0^{-1} \bm{U}\bm{U}^T}.
\end{equation*}
We end up the negative of the Fisher information operator, as computed in \eqref{eq:G_def}.

\begin{equation*}
  \frac{d^2}{dt^2}\bigg|_{t=0} L(t \bm{h}) = - \frac{1}{2}
  \tr \curly{\squared{\bm{\Sigma}_0^{-1}D_{\bm{h}} \bm{\Sigma}([\bm{W}_0, \sigma_0^2])}^2}
  = - G_{[\bm{W}_0, \sigma_0^2]}(\bm{h}, \bm{h}).
\end{equation*}
This completes the proof.
\end{proof}

\begin{definition}\label{def:G_inverse}
We define the \textbf{inverse Fisher information operator} $G_{[\bm{W}, \sigma^2]}^{-1}$
as the inverse of the flat version of the Fisher information bilinear form.
More concretely, instead of treating $G_{[\bm{W}, \sigma^2]}$ as a bilinear form on $T_{[\bm{W}, \sigma^2]}M$,
we treat it as a map $G_{[\bm{W}, \sigma^2]}: T_{[\bm{W}, \sigma^2]}M \to T^*_{[\bm{W}, \sigma^2]}M$
by setting $G_{[\bm{W}, \sigma^2]}(\bm{v}) = G_{[\bm{W}, \sigma^2]}(\bm{v}, \cdot)$
and we define $G_{[\bm{W}, \sigma^2]}^{-1}: T^*_{[\bm{W}, \sigma^2]}M \to T_{[\bm{W}, \sigma^2]}M$
to be the inverse.
\end{definition}

\begin{theorem}\label{thm:asc_linearity}
The maximum likelihood estimator is asymptotically linear at $[\bm{W}_0, \sigma_0^2]$
with the influence function $\text{IF}: \mathbb{R}^p \to T_{[\bm{W}_0, \sigma_0^2]}M$
given by 

\begin{equation}\label{eq:influence_function}
  \text{IF}(\bm{x}) = G_{[\bm{W}_0, \sigma_0^2]}^{-1}(s(\bm{x}; [\bm{W}_0, \sigma_0^2])).
\end{equation}
As described in \cite{Sun_Lin_Liu_2026} Lemma III.3, this means
\begin{equation}\label{eq:asc_linearity}
  \sqrt{n} \text{Exp}^{-1}_{[\bm{W}_0, \sigma_0^2]}\bracket{\squared{\hat{\bm{W}}_n,\hat{\sigma}_n^2}}
  - \frac{1}{\sqrt{n}} \sum_{i=1}^n \text{IF}(\bm{X}_i) \to 0
\end{equation}
in probability as $n \to \infty$ when $(\bm{X}_1, \ldots, \bm{X}_n)$
follows the true distribution $P^n_{[\bm{W}_0, \sigma_0^2]}$.
Also, $\text{IF}(\bm{X})$ has mean $\bm{0}$ and is square integrable if $\bm{X}$
follows the true distribution $N_p(\bm{0}, \bm{\Sigma}_0)$.

In particular, 
once we fix a $g_{[\bm{W}_0, \sigma_0^2]}$ orthonormal basis $\bm{e}_1, \ldots, \bm{e}_d$
and we identify $T_{[\bm{W}_0, \sigma_0^2]}M$ canonically as $\mathbb{R}^d$
(in concrete terms, identify $\sum_{j=1}^d a_j \bm{e}_j$ to the point $(a_1, \ldots, a_d) \in \mathbb{R}^d$),
then
$\sqrt{n} \text{Exp}^{-1}_{[\bm{W}_0, \sigma_0^2]}\bracket{\squared{\hat{\bm{W}}_n,\hat{\sigma}_n^2}}$
converges in distribution to $N_d(\bm{0}, \bm{G}^{-1})$ as $n \to \infty$ when $(\bm{X}_1, \ldots, \bm{X}_n)$
follows the true distribution $P^n_{[\bm{W}_0, \sigma_0^2]}$,
where $\bm{G}$ is the $d \times d$ 
matrix defined by $\bm{G}_{ij} = G_{[\bm{W}_0, \sigma_0^2]}(\bm{e}_i, \bm{e}_j)$.
\end{theorem}

\begin{proof}
Fix a $g_{[\bm{W}_0, \sigma_0^2]}$ orthonormal basis $\bm{e}_1, \ldots, \bm{e}_d$
of $T_{[\bm{W}_0, \sigma_0^2]}M$.
Define the matrix $\bm{G} \in \mathbb{R}^{d \times d}$ as in the statement of lemma 
and the vector $\nabla l(\bm{x}; [\bm{W}_0, \sigma_0^2]) \in \mathbb{R}^d$
by setting $(\nabla l(\bm{x}; [\bm{W}_0, \sigma_0^2]))_j = s(\bm{x}; [\bm{W}_0, \sigma_0^2])(\bm{e}_j)$,
for all $j = 1, \ldots, d$.
Identifying 
$\sqrt{n} \text{Exp}^{-1}_{[\bm{W}_0, \sigma_0^2]}\bracket{\squared{\hat{\bm{W}}_n,\hat{\sigma}_n^2}}
\in T_{[\bm{W}_0, \sigma_0^2]}M$ to a point in $\mathbb{R}^d$
through this basis,
we immediately have, by combining the two previous lemmas and Theorem 5.23 in \cite{van_der_Vaart_1998},
that

\begin{equation}\label{eq:asc_linearity_in_proof}
  \sqrt{n} \text{Exp}^{-1}_{[\bm{W}_0, \sigma_0^2]}\bracket{\squared{\hat{\bm{W}}_n,\hat{\sigma}_n^2}}
  - \frac{1}{\sqrt{n}} \sum_{i=1}^n \bm{G}^{-1} \nabla l(\bm{X}_i; [\bm{W}_0, \sigma_0^2]) \to 0
\end{equation}
in probability as $n \to \infty$.
In particular, $\sqrt{n} \text{Exp}^{-1}_{[\bm{W}_0, \sigma_0^2]}\bracket{\squared{\hat{\bm{W}}_n,\hat{\sigma}_n^2}}$
converges in distribution to $N_d(\bm{0}, \bm{G}^{-1})$ as $n \to \infty$.
Also, if $\bm{X}$ follows the true distribution $N_p(\bm{0}, \bm{\Sigma}_0)$,
then $\mathbb{E}[\nabla l(\bm{X}; [\bm{W}_0, \sigma_0^2])] = \bm{0}$
and thus $\mathbb{E}[\bm{G}^{-1}\nabla l(\bm{X}; [\bm{W}_0, \sigma_0^2])] = \bm{0}$.
Given that $\nabla l(\bm{X}; [\bm{W}_0, \sigma_0^2])$ has its second moment,
so does $\bm{G}^{-1}\nabla l(\bm{X}; [\bm{W}_0, \sigma_0^2])$.

We remark that 
\eqref{eq:asc_linearity}
is simply \eqref{eq:asc_linearity_in_proof}
rephrased in a way 
using only vectors in $T_{[\bm{W}_0, \sigma_0^2]}M$ defined without
$\bm{e}_1, \ldots, \bm{e}_d$
and hence without the need of identifying $T_{[\bm{W}_0, \sigma_0^2]}M$ with $\mathbb{R}^d$.
Let us for simplicity define $\bm{y} := \bm{G}^{-1} \nabla l(\bm{x}; [\bm{W}_0, \sigma_0^2]) \in \mathbb{R}^d$.
Put $\bm{v} \in T_{[\bm{W}_0, \sigma_0^2]}M$ by setting $\bm{v} = \sum_{j=1}^d \bm{y}_j \bm{e}_j \in T_{[\bm{W}_0, \sigma_0^2]}M$.
We claim that $\bm{v}$ is the unique vector such that, for all $\bm{z} \in T_{[\bm{W}_0, \sigma_0^2]}M$:

\begin{equation}\label{eq:transition_to_G_inverse}
  G_{[\bm{W}_0, \sigma_0^2]}(\bm{v}, \bm{z}) = s(\bm{x}; [\bm{W}_0, \sigma_0^2])(\bm{z}).
\end{equation}
The uniqueness of $\bm{v}$ is immediate by the positive definiteness of $G_{[\bm{W}_0, \sigma_0^2]}$.
Given that $(\bm{G}\bm{y})_j = (\nabla l(\bm{x}; [\bm{W}_0, \sigma_0^2]))_j$
for all $j = 1, \ldots, d$,
we have, for all $j = 1, \ldots, d$:
\begin{equation*}
  G_{[\bm{W}_0, \sigma_0^2]}(\bm{v}, \bm{e}_j) = 
  \sum_{i=1}^d \bm{y}_i G_{[\bm{W}_0, \sigma_0^2]}(\bm{e}_i, \bm{e}_j)
  = \sum_{i=1}^d \bm{y}_i \bm{G}_{ij} = s(\bm{x}; [\bm{W}_0, \sigma_0^2])(\bm{e}_j).
\end{equation*}
The claim is therefore proved.
By looking at Definition~\ref{def:G_inverse}, 
\eqref{eq:transition_to_G_inverse}
reads $G_{[\bm{W}_0, \sigma_0^2]}(\bm{v})=s(\bm{x}; [\bm{W}_0, \sigma_0^2])$
and hence $\bm{v} = G_{[\bm{W}_0, \sigma_0^2]}^{-1}(s(\bm{x}; [\bm{W}_0, \sigma_0^2]))$,
giving \eqref{eq:asc_linearity}.
The results about the integral of $\text{IF}(\bm{X})$
translates to here as well.
This completes the proof.
\end{proof}

\begin{theorem}
We have that

\begin{equation}\label{eq:covariance_fit}
  \sqrt{n}\bracket{\hat{\bm{\Sigma}}_n - \bm{\Sigma}_0}
  - \frac{1}{\sqrt{n}}\sum_{i=1}^n D_{\text{IF}(\bm{X}_i)}\bm{\Sigma}(\squared{\bm{W}_0, \sigma_0^2}) \to \bm{0}
\end{equation}
in probability as $n \to \infty$ when 
$(\bm{X}_1, \ldots, \bm{X}_n)$
follows the true distribution $P^n_{[\bm{W}_0, \sigma_0^2]}$.

As in Theorem~\ref{thm:asc_linearity},
let $g_{[\bm{W}_0, \sigma_0^2]}$-orthonormal basis $\bm{e}_1, \ldots, \bm{e}_d$
be fixed and $\bm{G}$ be the $d \times d$ matrix corresponding to $G_{[\bm{W}_0, \sigma_0^2]}$.
Define a function $D: \mathbb{R}^d \to \mathbb{R}^{p \times p}$
by setting, for all $\bm{z} = (\bm{z}_1, \ldots, \bm{z}_d)$:

\begin{equation*}
  D(\bm{z}) = D_{\sum_{i=1}^d \bm{z}_i \bm{e}_i}\bm{\Sigma}(\squared{\bm{W}_0, \sigma_0^2}).
\end{equation*}
Then, 

\begin{equation*}
  \sqrt{n}\bracket{\hat{\bm{\Sigma}}_n - \bm{\Sigma}_0} \to D(\bm{Z})
\end{equation*}
in distribution as $n \to \infty$ 
when 
$(\bm{X}_1, \ldots, \bm{X}_n)$
follows the true distribution $P^n_{[\bm{W}_0, \sigma_0^2]}$,
where $\bm{Z}$ is a random variable following $N_d(\bm{0}, \bm{G}^{-1})$
distribution.
\end{theorem}

\begin{proof}
Suppose that a given tangent vector $\bm{h} \in \mathcal{E}_{[\bm{W}_0, \sigma_0^2]}$
lifts horizontally to $(\bm{U}, b)$ at $(\bm{W}_0, \sigma_0^2)$.
\eqref{eq:Sigma_derivative} shows that 

\begin{equation*}
  \bm{\Sigma}\bracket{\text{Exp}_{[\bm{W}_0, \sigma_0^2]}(\bm{h})}
  = (\bm{W}_0+\bm{U})(\bm{W}_0+\bm{U})^T + (\sigma_0^2+b)\bm{I}_p
  = \bm{\Sigma}_0 + \bm{U}\bm{U}^T + D_{\bm{h}} \bm{\Sigma}([\bm{W}_0, \sigma_0^2]).
\end{equation*}
For simplicity, put 

\begin{equation*}
  \bm{h}_n := \text{Exp}^{-1}_{[\bm{W}_0, \sigma_0^2]}\bracket{\squared{\hat{\bm{W}}_n,\hat{\sigma}_n^2}}.
\end{equation*}
If $\bm{h}_n$
lifts horizontally to $(\bm{U}_n, b_n)$ at $(\bm{W}_0, \sigma_0^2)$,
then 

\begin{equation*}
  \sqrt{n}\bracket{\hat{\bm{\Sigma}}_n - \bm{\Sigma}_0}
  = \sqrt{n}\bm{U}_n\bm{U}_n^T + \sqrt{n} D_{\bm{h}_n} \bm{\Sigma}([\bm{W}_0, \sigma_0^2]).
\end{equation*}
Since $\|\bm{U}_n\bm{U}_n^T\|_F \leqslant	\|\bm{U}_n\|_F^2 \leqslant \abs{\bm{h}_n}^2$,
the asymptotic normality of $\sqrt{n}\bm{h}_n$ proved in 
Theorem~\ref{thm:asc_linearity} implies 
that $\sqrt{n}\bm{U}_n\bm{U}_n^T \to \bm{0}$ 
in probability as $n \to \infty$ when $(\bm{X}_1, \ldots, \bm{X}_n)$
follows the true distribution $P^n_{[\bm{W}_0, \sigma_0^2]}$.
As of the other term,
given that the directional derivative is a linear map with respect to 
$\bm{h}$, it follows from \eqref{eq:asc_linearity} that 

\begin{equation*}
  \sqrt{n} D_{\bm{h}_n} \bm{\Sigma}([\bm{W}_0, \sigma_0^2])
  - \frac{1}{\sqrt{n}} \sum_{i=1}^n D_{\text{IF}(\bm{X}_i)}\bm{\Sigma}(\squared{\bm{W}_0, \sigma_0^2}) \to 0
\end{equation*}
in probability as $n \to \infty$ when $(\bm{X}_1, \ldots, \bm{X}_n)$
follows the true distribution $P^n_{[\bm{W}_0, \sigma_0^2]}$.
We have therefore proved \eqref{eq:covariance_fit}.
The rest of the conclusions follow immediately.
\end{proof}

\subsection{Local Regularity}

We now investigate the behaviour of the maximum likelihood 
estimator when the underlying data has a distribution 
slightly perturbed from the true parameter
and eventually show that the maximum likelihood estimator is locally regular
in the sense described in \cite{Sun_Lin_Liu_2026}.
Given that calculations regarding components of a point in $M$ 
does not matter that much in this section,
we shall proceed with the notations already introduced in Proposition~\ref{thm:LAN}.
In concrete terms, 
let $\squared{\hat{\theta}_n}=[\hat{\bm{W}}_n, \hat{\sigma}_n^2]$ be the maximum likelihood estimator 
and, for any $\bm{u} \in \mathcal{E}_{[\theta_0]}$, let

\begin{equation*}
  \Gamma_{\bm{u}} := \Pi^{[\theta_0]}_{\text{Exp}_{[\theta_0]}(\bm{u})} :
  T_{\text{Exp}_{[\theta_0]}(\bm{u})}M  \to T_{[\theta_0]}M
\end{equation*}
be the parallel transport map considered in 
Proposition~\ref{thm:parallel_transport}.
Given $\bm{h} \in T_{[\theta_0]}M$,
our goal of the subsection is to show 
the asymptotic normality of

\begin{equation*}
  \sqrt{n} \Gamma_{\frac{\bm{h}}{\sqrt{n}}} \bracket{\text{Exp}^{-1}_{
    \squared{\theta_{n, \bm{h}}}}\bracket{\squared{\hat{\theta}_n}}}
\end{equation*}
as $n \to \infty$ when $(\bm{X}_1, \ldots, \bm{X}_n)$ 
follows $P^n_{\squared{\theta_{n,\bm{h}}}}$.
We start with technical intermediate results.

\begin{lemma}\label{thm:R_u_w_bound}
Suppose that $K$ is a compact 
convex neighborhood of $(\bm{0}, \bm{0})$ in 
$\mathbb{R}^d \times \mathbb{R}^d$,
$R: \mathbb{R}^d \times \mathbb{R}^d \to \mathbb{R}^d$ is a smooth function, 
$R((\bm{u}, \bm{u})) = \bm{0}$ for every $\bm{u} \in \mathbb{R}^d$,
and $R(\bm{0}, \bm{w}) = \bm{0}$ for every $\bm{w} \in \mathbb{R}^d$.
Then, there exists a constant $C \in (0, \infty)$ such that, 
for all $(\bm{u}, \bm{w})\in K$:

\begin{equation*}
  \abs{R(\bm{u}, \bm{w})} \leqslant C \abs{\bm{u}}\abs{\bm{w}-\bm{u}}.
\end{equation*}
\end{lemma}

\begin{proof}
By the fundamental theorem of calculus, we have

\begin{equation*}
  \begin{split}
  & \tab[0.4cm] R(\bm{u}, \bm{w}) = 
  \int_0^1 \frac{\partial R}{\partial \bm{w}}(\bm{u}, \bm{u} + t(\bm{w}-\bm{u}))(\bm{w}-\bm{u}) \mathrm{d} t
  \\ & = \int_0^1 \int_0^1 \frac{\mathrm{d}}{\mathrm{d}s}
  \squared{\frac{\partial R}{\partial \bm{w}}(s\bm{u}, \bm{u} + t(\bm{w}-\bm{u}))(\bm{w}-\bm{u})} \mathrm{d} s \mathrm{d} t
  = \int_0^1 \int_0^1 \frac{\partial }{\partial \bm{u}} \squared{\frac{\partial R}{\partial \bm{w}}
  (s\bm{u}, \bm{u} + t(\bm{w}-\bm{u}))(\bm{w}-\bm{u})}\bm{u}  \mathrm{d} s \mathrm{d} t,
  \end{split}
\end{equation*}
where we have used the fact that $\frac{\partial R}{\partial \bm{w}} R (0, \bm{u} + t(\bm{w}-\bm{u})) \equiv 0$
because $R (0, \bm{u} + t(\bm{w}-\bm{u})) \equiv 0$.
With 

\begin{equation*}
  C := \max_{(\bm{u}, \bm{w}) \in K } \abs{\frac{\partial }{\partial \bm{u}} \frac{\partial R}{\partial \bm{w}}
  (\bm{u}, \bm{w})},
\end{equation*}
which is finite because $K$ is compact,
we have, for all $(\bm{u}, \bm{w})\in K$:

\begin{equation*}
  \abs{R(\bm{u}, \bm{w})} \leqslant C \int_0^1 \int_0^1 \abs{\bm{u}}\abs{\bm{w}-\bm{u}}  \mathrm{d} s \mathrm{d} t
  = C \abs{\bm{u}}\abs{\bm{w}-\bm{u}}.
\end{equation*}
The proof is complete.
\end{proof}

\begin{lemma}\label{thm:perturbed_lemma}
Let $T_n = T_n(\bm{X}_1, \ldots, \bm{X}_n) \in M$ be an estimator 
such that, when $(\bm{X}_1, \ldots, \bm{X}_n)$ follows $P^n_{[\theta_0]}$,
the sequence
$\sqrt{n}\text{Exp}_{[\theta_0]}^{-1}(T_n) \in T_{\theta_0}M \cong \mathbb{R}^d$ is 
bounded in probability as $n \to \infty$.
(In symbols, it is $O_P(1)$.)
Then, for each $\bm{h} \in T_{[\theta_0]}M$,

\begin{equation*}
  \sqrt{n} \Gamma_{\frac{\bm{h}}{\sqrt{n}}}
  \bracket{\text{Exp}^{-1}_{\squared{\theta_{n,\bm{h}}}}(T_n)}
  - \sqrt{n} \text{Exp}_{[\theta_0]}^{-1}(T_n) 
  + \bm{h}  \to \bm{0}
\end{equation*}
in probability as $n \to \infty$ 
when $(\bm{X}_1, \ldots, \bm{X}_n)$ follows $P^n_{\squared{\theta_{n,\bm{h}}}}$.
\end{lemma}

\begin{proof}
Fix a constant $r \in \bracket{0, 
\frac{\min(\sigma_{\min}(\bm{W}_0), \sigma_0^2)}{3}}$
so that the ball $B_r$ centred at $\bm{0} \in T_{[\theta_0]}M$
with radius $r$ is compactly contained in $\mathcal{E}_{[\theta_0]}$
and, for every $(\bm{u}, \bm{w}) \in B_r \times B_r$,
$\text{Exp}_{[\theta_0]}(\bm{w}) \in 
U_{\text{Exp}_{[\theta_0]}(\bm{u})}$.
Define $F : B_r \times B_r \to T_{[\theta_0]}M$ by setting,
for all $(\bm{u}, \bm{w}) \in B_r \times B_r$,
that 

\begin{equation*}
  F(\bm{u}, \bm{w}) = \Gamma_{\bm{u}}
  \bracket{\text{Exp}^{-1}_{\text{Exp}_{[\theta_0]}(\bm{u})}
  \bracket{\text{Exp}_{[\theta_0]}(\bm{w})}}.
\end{equation*}
Then $F$ is a smooth function because 
the map from $(\bm{u}, \bm{z})$
to $\Gamma_{\bm{u}}(\bm{z})$ is smooth (as shown with effort in Proposition~\ref{thm:parallel_transport})
and the map from $(\bm{u}, \bm{w})$
to $\text{Exp}^{-1}_{\text{Exp}_{[\theta_0]}(\bm{u})}
  \bracket{\text{Exp}_{[\theta_0]}(\bm{w})}$
is smooth.
As a consequence,
the function 

\begin{equation*}
  R(\bm{u}, \bm{w}) = F(\bm{u}, \bm{w}) - (\bm{w}-\bm{u})
\end{equation*}
satisfies the hypothesis of Lemma~\ref{thm:R_u_w_bound}
and thus 

\begin{equation*}
  \begin{split}
  & \tab[0.4cm]  \abs{\sqrt{n} \Gamma_{\frac{\bm{h}}{\sqrt{n}}}
  \bracket{\text{Exp}^{-1}_{\squared{\theta_{n,\bm{h}}}}(T_n)}
  - \sqrt{n} \text{Exp}_{[\theta_0]}^{-1}(T_n) 
  + \bm{h}} =
  \sqrt{n}\abs{F \bracket{\frac{\bm{h}}{\sqrt{n}}, \text{Exp}^{-1}_{[\theta_0]}(T_n)} - 
  \bracket{\text{Exp}^{-1}_{[\theta_0]}(T_n)-\frac{\bm{h}}{\sqrt{n}}}}
  \\ & \leqslant C \abs{\frac{\bm{h}}{\sqrt{n}}}
  \abs{\sqrt{n}\text{Exp}^{-1}_{[\theta_0]}(T_n)-\bm{h}}.
  \end{split}
\end{equation*}
By assumption, $\sqrt{n}\text{Exp}^{-1}_{[\theta_0]}(T_n)$
is bounded in probability 
when $(\bm{X}_1, \ldots, \bm{X}_n)$ follows $P^n_{[\theta_0]}$.
Given that we have proved in Proposition~\ref{thm:LAN}
that $P^n_{[\theta_0]}$ and $P^n_{\squared{\theta_{n,\bm{h}}}}$ are mutually 
contiguous, $\sqrt{n}\text{Exp}^{-1}_{[\theta_0]}(T_n)$ is also 
bounded in probability when $(\bm{X}_1, \ldots, \bm{X}_n)$ follows 
$P^n_{\squared{\theta_{n,\bm{h}}}}$.
The proof is therefore complete.
\end{proof}

\begin{theorem}
The maximum likelihood estimator is locally regular.
With
any $g_{[\theta_0]}$-orthonormal basis 
$\bm{e}_1, \ldots, \bm{e}_d$ of $T_{[\theta_0]}M$ fixed
and the $d \times d$ matrix $\bm{G}$ defined just as in Theorem~\ref{thm:asc_linearity},
we have

\begin{equation*}
  \sqrt{n} \Gamma_{\frac{\bm{h}}{\sqrt{n}}} \bracket{\text{Exp}^{-1}_{
    \squared{\theta_{n, \bm{h}}}}\bracket{\squared{\hat{\theta}_n}}}
    \to N_d \bracket{\bm{0}, \bm{G}^{-1}}
\end{equation*}
in distribution 
as $n \to \infty$ when $(\bm{X}_1, \ldots, \bm{X}_n)$ 
follows $P^n_{\squared{\theta_{n,\bm{h}}}}$.
\end{theorem}

\begin{proof}
Put 

\begin{equation*}
  \Lambda_n(\bm{h}) = \sum_{i=1}^n \log \frac{p\bracket{\bm{X}_i; \squared{\theta_{n, \bm{h}}}}}{p(\bm{X}_i; [\theta_0])}.
\end{equation*}
We claim that 
\begin{equation}\label{eq:joint}
    \begin{bmatrix}
      \sqrt{n}\text{Exp}^{-1}_{\squared{\theta_0}}\bracket{\squared{\hat{\theta}_n}} \\
      \Lambda_n(\bm{h})
    \end{bmatrix}
  \to
  N_{d+1}\bracket{
  \begin{bmatrix} 
    \bm{0}\\ -\frac{1}{2}\bm{h}^T\bm{G}\bm{h}
  \end{bmatrix},
  \begin{bmatrix} \bm{G}^{-1} & \bm{h} \\
  \bm{h}^T & \bm{h}^T\bm{G}\bm{h}
\end{bmatrix}}
\end{equation}
in distribution as $n \to \infty$ when 
$(\bm{X}_1, \ldots, \bm{X}_n)$ follows $P^n_{\theta_0}$.
The the asymptotic mean and the diagonal part of the asymptotical
covariance matrix
are already justified by 
\eqref{eq:LAN_weak_convergence}
and Theorem~\ref{thm:asc_linearity}.
Next, by \eqref{eq:asc_linearity_in_proof}
and \eqref{eq:LAN_prob_convergence},
we see that the asymptotic covariance term 
equals to $\text{Cov}(\bm{G}^{-1}\Delta,\Delta^T\bm{h})$
with $\Delta$ being a random variable 
with distribution $N_d \bracket{\bm{0}, \bm{G}}$,
which calculates to $\bm{G}^{-1}\bm{G}\bm{h} = \bm{h}$.

Next, by Le Cam's third lemma as in Example 6.7 in \cite{van_der_Vaart_1998}, \eqref{eq:joint}
implies that $\sqrt{n}\text{Exp}_{\squared{\theta_0}}\bracket{\squared{\hat{\theta}_n}}$
converges to $N_d(\bm{h}, \bm{G}^{-1})$ in distribution as $n \to \infty$
when $(\bm{X}_1, \ldots, \bm{X}_n)$ 
follows $P^n_{\squared{\theta_{n,\bm{h}}}}$.
Applying Lemma~\ref{thm:perturbed_lemma} with $T_n = \squared{\hat{\theta}_n}$
will finish the proof.

\end{proof}

\chapter{Next Steps}

\section{High Dimensional Case}

\section{Next Steps in Low Dimensions}

\section{Simulations}

\bibliographystyle{plainnat}
\def\bibname{References} % Nom obligatoire de la section des références.
                                              % On utilise \'e car le é cause des problèmes
                                              % dans la table des matière
%% ENGLISH
 %\def\bibname{References}
\bibliography{reference.bib}     % La base de données contenant des entrées bibliographiques.
                                    % Seules celles référencées dans le texte seront ajoutées
                                    % \`a la bibliographie.

\end{CJK*}
\end{document}
